\documentclass[11pt]{article}% Shared preamble for the rigor documents (Lamport-style structured proofs).  \input this file after \documentclass.
\usepackage[T1]{fontenc}\usepackage{lmodern}\usepackage{microtype}
\usepackage[margin=1in]{geometry}
\usepackage{amsmath,amssymb,amsthm,mathtools}
\usepackage{xcolor}\usepackage{enumitem}\usepackage{booktabs}\usepackage{graphicx}\usepackage{hyperref}
\usepackage[most]{tcolorbox}
\definecolor{navy}{HTML}{17365D}\definecolor{teal}{HTML}{117A8B}\definecolor{warn}{HTML}{A33A2B}\definecolor{okgreen}{HTML}{2E7D4F}
\hypersetup{colorlinks=true,linkcolor=navy,citecolor=teal,urlcolor=teal}
\theoremstyle{plain}
\newtheorem{theorem}{Theorem}[section]\newtheorem{lemma}[theorem]{Lemma}\newtheorem{proposition}[theorem]{Proposition}\newtheorem{corollary}[theorem]{Corollary}
\theoremstyle{definition}\newtheorem{definition}[theorem]{Definition}\newtheorem{remark}[theorem]{Remark}\newtheorem{claim}[theorem]{Claim}\newtheorem{issue}[theorem]{Issue}
% ---------- Lamport structured proofs ----------
% \lstep{level}{number}{assertion}   -> indented "<level>number. assertion"
% \lproof                             -> "PROOF:" line (start of the sub-proof of the preceding step)
% \lby{justification}                 -> "BY ..." leaf justification for the preceding step
% \lqed{level}{number}                -> QED step of a sub-proof
% keywords: \lassume \lprove \lsuffices \lcase \ldefine \lpick \lnew
\newlength{\lindent}\setlength{\lindent}{1.7em}
\newcommand{\lstep}[3]{\par\smallskip\noindent\hspace*{\dimexpr\numexpr#1-1\relax\lindent\relax}{\color{navy}$\langle#1\rangle#2$.}\ #3}
\newcommand{\lproof}{\par\noindent\hspace*{\lindent}{\scshape Proof:}}
\newcommand{\lby}[1]{\par\noindent\hspace*{\lindent}{\scshape By}\ #1.}
\newcommand{\lqed}[2]{\par\smallskip\noindent\hspace*{\dimexpr\numexpr#1-1\relax\lindent\relax}{\color{navy}$\langle#1\rangle#2$.}\ {\scshape Q.E.D.}}
\newcommand{\lassume}{{\scshape Assume}\ }\newcommand{\lprove}{{\scshape Prove}\ }\newcommand{\lsuffices}{{\scshape Suffices}\ }
\newcommand{\lcase}{{\scshape Case}\ }\newcommand{\ldefine}{{\scshape Define}\ }\newcommand{\lpick}{{\scshape Pick}\ }\newcommand{\lnew}{{\scshape New}\ }
% ---------- verdict boxes ----------
\newtcolorbox{verdictok}{colback=okgreen!8,colframe=okgreen,title={\bfseries Verdict: verified},fonttitle=\small}
\newtcolorbox{verdictgap}{colback=orange!10,colframe=orange!80!black,title={\bfseries Verdict: gap / caveat},fonttitle=\small}
\newtcolorbox{verdictbreak}{colback=warn!10,colframe=warn,title={\bfseries Verdict: proof-breaking issue},fonttitle=\small}
\newtcolorbox{citedbox}[1]{colback=navy!5,colframe=navy!60,title={\bfseries Cited result: #1},fonttitle=\small,breakable}
\setlength{\parindent}{0pt}\setlength{\parskip}{4pt}\allowdisplaybreaks
\newcommand{\cA}{\mathcal A}\newcommand{\cF}{\mathcal F}\newcommand{\cM}{\mathfrak M}\newcommand{\cN}{\mathfrak N}

% extra calligraphic shortcuts (\providecommand: no clash if a document defines its own)
\providecommand{\cB}{\mathcal B}\providecommand{\cC}{\mathcal C}\providecommand{\cD}{\mathcal D}\providecommand{\cE}{\mathcal E}\providecommand{\cG}{\mathcal G}\providecommand{\cH}{\mathcal H}\providecommand{\cI}{\mathcal I}\providecommand{\cJ}{\mathcal J}\providecommand{\cK}{\mathcal K}\providecommand{\cL}{\mathcal L}\providecommand{\cO}{\mathcal O}\providecommand{\cP}{\mathcal P}\providecommand{\cQ}{\mathcal Q}\providecommand{\cR}{\mathcal R}\providecommand{\cS}{\mathcal S}\providecommand{\cT}{\mathcal T}\providecommand{\cU}{\mathcal U}\providecommand{\cV}{\mathcal V}\providecommand{\cW}{\mathcal W}\providecommand{\cX}{\mathcal X}\providecommand{\cY}{\mathcal Y}\providecommand{\cZ}{\mathcal Z}
\newcommand{\Fid}{\operatorname{F}}\newcommand{\Tr}{\operatorname{Tr}}\newcommand{\eps}{\varepsilon}\newcommand{\dd}{\,\mathrm d}

\newcommand{\Pp}{P}\newcommand{\Pperp}{P^{\perp}}
\newcommand{\Wt}{\widetilde W}\newcommand{\kt}{\widetilde k}\newcommand{\Om}{\Omega}\newcommand{\Gm}{\Gamma}
\newcommand{\Stwo}{\mathfrak S_2}\newcommand{\Sone}{\mathfrak S_1}
\newcommand{\norm}[1]{\lVert #1\rVert}\newcommand{\abs}[1]{\lvert #1\rvert}
\newcommand{\ip}[2]{\langle #1,#2\rangle}\newcommand{\one}{\mathbf 1}
\providecommand{\ket}[1]{|#1\rangle}\providecommand{\bra}[1]{\langle #1|}
\newcommand{\Ext}{\mathrm{Ext}}\newcommand{\Nvis}{\mathcal N}
\newcommand{\Erec}{E_{\mathrm{rec}}}\newcommand{\Gs}{\mathcal G}
\DeclareMathOperator{\dist}{dist}\DeclareMathOperator{\Var}{Var}
\DeclareMathOperator{\ran}{ran}\DeclareMathOperator{\Ad}{Ad}
\title{The upper half of the exact-fidelity question:\\
\large the exterior directions exhaust the invisible space, and
$\Erec\le(1-\theta)f_2z^2(1+o(1))$}
\author{Agent E1 (Phase 4, analytic track)}\date{2026-09-15}
\begin{document}\maketitle
\begin{abstract}
Let $P$ be the Hardy projection of $L^2(\mathbb R)$ and $E_I$ the projection onto $L^2(I)$,
$I$ an interval.  Inside the real Hilbert space $\Stwo(PH,\Pperp H)$ let
$\Nvis=\{\Pperp\ell P:\ \ell=\ell^*=E_I\ell E_I\}$ be the \emph{visible} space of first-order
recovery defects and $\Ext=\{\Pperp Z_eP:\ Z_e^*=-Z_e=E_{I^c}Z_eE_{I^c}\}$ the space of
\emph{exterior} directions.  We prove $\overline{\Ext}=\Nvis^{\perp}$
(Theorem~\ref{thm:Q}): the exterior extensions of a channel exhaust exactly the directions
that the interior cannot see (the visible space being generated by finite-rank self-adjoint operators localised in $I$, the exterior one by finite-rank anti-self-adjoint operators localised in $I^c$).  Two independent proofs are given, one by an invariant two-projection computation for the pair $(P,E_I)$ (self-contained, one-particle; no normal form is needed), one by
Tomita--Takesaki duality of the standard subspaces of $\cF(I)$ and $\cF(I)'$ in the
one-particle--one-hole sector.  Consequences: the Uhlmann bound optimised over exterior
extensions equals the Legendre value $g_Q$ of the first-order defect \emph{for every
isometric channel} (Theorem~\ref{thm:reduction}), which is the missing upper half (UH) of the
exact-fidelity question; with the one-parameter-group remainder estimate this upgrades to a
statement about the exact fidelity (Theorem~\ref{thm:UH}), and hence $\Erec\le(1-\theta_{\rm fr})f_2z^2(1+o(1))$
(Theorem~\ref{thm:main}), where $\theta_{\rm fr}$ is the isometric-orbit gain over finite-rank
rotations $\Gs$ with $v_\Gs\ne0$, and $\theta_{\rm fr}>0$ iff $E_DuE_D\ne0$.  With the numerical
$\theta_{\rm fr}\ge0.38$ (Phase~5, F3: $\theta=0.384\pm0.003$) this says that the zero-collar
compression is \emph{not} optimal in exact fidelity.
\end{abstract}
\tableofcontents
\section{Setting, and the exact Uhlmann bound for an arbitrary extension}\label{sec:setting}

Conventions are those of \texttt{rigor/uhlmann\_upper\_bound.tex} (``PA'') and
\texttt{rigor/sdp\_dual\_certificate.tex} (``S10''); they are recalled only as far as used.

\paragraph{(E1) One-particle data.}
$H:=L^2(\mathbb R)$ with $\ip fg=\int\overline fg$; $P$ is the one-particle vacuum (Hardy)
projection of the chiral free fermion and $\Pperp:=\one-P$ (PA (C1); the sign convention
$P=P_+$ or $P=P_-$ is irrelevant below and is fixed once by writing $P$).  For a Borel set
$S$, $E_S$ is multiplication by $\one_S$ and $\mathfrak h_S:=L^2(S)=E_SH$.  Geometry (S10 (C1)):
$A=(-a,0)$, $B=(0,b)$, $C=(b,L)$, $D=BC=(0,L)$, $I=ABC=(-a,L)$, $I^c=\mathbb R\setminus
\overline I$, cross ratio $z=ac/(b(a+b+c))$, $f_2=1/(12\pi^2)$ per chirality.  All
Hilbert--Schmidt classes are denoted $\Stwo$, and
\begin{equation}\label{eq:realHS}
 \cH:=\Stwo(PH,\Pperp H)=\{\Pperp Y P:\ Y\in\Stwo\},\qquad
 \ip XY_r:=\Re\Tr(X^*Y),
\end{equation}
regarded as a \emph{real} Hilbert space with the inner product $\ip\cdot\cdot_r$; ``$\perp$'',
``$\dist$'' and orthogonal projections always refer to $\ip\cdot\cdot_r$.

\paragraph{(E2) Fock data.}
$\cF(X):=\mathrm{CAR}(L^2(X))''$ in the Fock representation with vacuum $\Om$ and
$\omega(a^*(f)a(g))=\ip g{Pf}$ (PA (C2)).  $\Om$ is cyclic and separating for $\cM:=\cF(I)$
(Reeh--Schlieder).  For a unitary $U$ of $H$ with $\Pperp UP\in\Stwo$, Shale--Stinespring
provides an implementer $\Gm(U)$ on Fock space, unique up to a phase, with
$\Gm(U)a(f)\Gm(U)^*=a(Uf)$ (PA Cor.~2.7).

\paragraph{(E3) Channels, defects, and the second-order functional.}
A recovery channel is a unital normal even CP map $\beta:\cA(D)\to\cA(B)$, and
$\omega^\beta:=\omega|_{\cA(AB)}\circ(\mathrm{id}_A\otimes\beta)$, so that
$\Erec:=\inf_\beta[-\log\mathrm F(\omega,\omega^\beta)]$ (\texttt{optimality\_all\_channels.tex},
``S12'' below, \S1).  $\beta$ is
\emph{isometric quasi-free} if $\beta(a(g))=a(Vg)$ for an isometry $V:\mathfrak h_D\to
\mathfrak h_B$ \emph{onto} $\mathfrak h_B$; we then write
\begin{equation}\label{eq:WI}
 W_I:=\one_{\mathfrak h_A}\oplus V:\ \mathfrak h_I\longrightarrow\mathfrak h_{AB}\subseteq\mathfrak h_I ,
\end{equation}
an isometry of $\mathfrak h_I$ into itself with range $\mathfrak h_{AB}$, and
$\omega^\beta(X)=\omega(\alpha_{W_I}(X))$ for $X\in\cM$, $\alpha_{W_I}(a(f)):=a(W_If)$.
The one-particle defect and the second-order functional are
\begin{equation}\label{eq:defect}
 \delta:=Q-Q_{\rm rec}=E_I\bigl(P-W_I^*PW_I\bigr)E_I,\qquad
 g_Q(\delta)=\tfrac18\sup_{t=t^*}\frac{(\Tr t\delta)^2}{\Tr(tQt(\one-Q))}
\end{equation}
(S10 (C2)--(C4), $Q=E_IPE_I$); $g_Q$ is the Legendre form of the Bures/SLD metric.

\begin{definition}[Admissible extensions]\label{def:adm}
Let $W_I$ be as in \eqref{eq:WI}.  Put
\begin{equation}\label{eq:adm}
 \mathcal E(W_I):=\bigl\{\Wt:\ \Wt\ \text{unitary on }H,\ \Wt E_I=W_IE_I,\ \Pperp\Wt P\in\Stwo\bigr\}.
\end{equation}
\end{definition}

\begin{lemma}[Structure of $\mathcal E(W_I)$]\label{lem:torsor}
If $\mathcal E(W_I)\ne\emptyset$ and $\Wt_0\in\mathcal E(W_I)$, then
$\mathcal E(W_I)=\{\Wt_0\widehat W':\ W'\ \text{unitary of }\mathfrak h_{I^c},\
\Pperp\Wt_0\widehat W'P\in\Stwo\}$, where $\widehat W':=\one_{\mathfrak h_I}\oplus W'$.
In particular the exterior part of an admissible extension is free: it is an arbitrary
unitary of $\mathfrak h_{I^c}$, composed with $\Wt_0$.
\end{lemma}
\lstep{1}{1}{If $\Wt\in\mathcal E(W_I)$ then $\Wt^*_0\Wt$ is unitary and acts as the identity on
 $\mathfrak h_I$.}
\lby{$\Wt E_I=W_IE_I=\Wt_0E_I$ gives $\Wt_0^*\Wt E_I=\Wt_0^*\Wt_0E_I=E_I$}
\lstep{1}{2}{A unitary $\widehat W$ of $H$ with $\widehat WE_I=E_I$ is of the form
 $\one_{\mathfrak h_I}\oplus W'$ with $W'$ unitary on $\mathfrak h_{I^c}$.}
\lby{$\widehat W$ fixes $\mathfrak h_I$ pointwise, hence preserves $\mathfrak h_I^\perp
 =\mathfrak h_{I^c}$ (a unitary maps the orthogonal complement of an invariant subspace on which
 it is onto, onto that complement)}
\lstep{1}{3}{Conversely every $\Wt_0\widehat W'$ with $\Pperp\Wt_0\widehat W'P\in\Stwo$ lies in
 $\mathcal E(W_I)$.}
\lby{$\Wt_0\widehat W'E_I=\Wt_0E_I=W_IE_I$ by $\langle1\rangle2$ read backwards}
\lqed{1}{4}\lby{$\langle1\rangle1$--$\langle1\rangle3$}

\begin{proposition}[Exact Uhlmann bound for an arbitrary admissible extension]\label{prop:uhl}
Let $\beta$ be an isometric quasi-free channel with one-particle isometry $W_I$
\eqref{eq:WI}, and let $\Wt\in\mathcal E(W_I)$.  Put $V:=\Pperp\Wt P$.  Then
\begin{equation}\label{eq:uhlbound}
 -\log\mathrm F\bigl(\omega|_\cM,\omega^\beta|_\cM\bigr)\ \le\
 -\tfrac12\log\det(\one-V^*V)\ \le\ -\tfrac12\log\bigl(\one-\norm V_2^2\bigr)
 \qquad(\norm V_2<1),
\end{equation}
and the middle term is $\tfrac12\norm V_2^2+O(\norm V_2^4)$.
\end{proposition}
\lstep{1}{1}{$\Psi:=\Gm(\Wt)^*\Om$ is a unit vector with $\ip\Psi{X\Psi}=\omega^\beta(X)$ for
 all $X\in\cM$.}
\lproof
\lstep{2}{1}{$\Gm(\Wt)X\Gm(\Wt)^*=\alpha_{W_I}(X)$ for $X$ in the $*$-algebra generated by
 $\{a(f):f\in\mathfrak h_I\}$.}
\lby{$\Gm(\Wt)a(f)\Gm(\Wt)^*=a(\Wt f)=a(W_If)$ for $f\in\mathfrak h_I$ by (E2) and
 $\Wt E_I=W_IE_I$ (Def.~\ref{def:adm}); then multiplicativity}
\lstep{2}{2}{$\ip\Psi{X\Psi}=\omega(\alpha_{W_I}(X))=\omega^\beta(X)$ on that $*$-algebra.}
\lby{$\ip\Psi{X\Psi}=\ip\Om{\Gm(\Wt)X\Gm(\Wt)^*\Om}$ and $\langle2\rangle1$; the second equality
 is the definition of $\omega^\beta$ in (E3) together with
 $\alpha_{W_I}|_{\cA(A)}=\mathrm{id}$ and $\alpha_{W_I}(a(g))=a(Vg)=\beta(a(g))$ for
 $g\in\mathfrak h_D$}
\lstep{2}{3}{Both sides of $\langle2\rangle2$ are $\sigma$-weakly continuous on $\cM$ and agree on a
 $\sigma$-weakly dense $*$-subalgebra, hence agree on $\cM$.}
\lby{the left side is a vector state; the right side is $\omega\circ\alpha_{W_I}$ with
 $\alpha_{W_I}$ normal (it is implemented by the unitary $\Gm(\Wt)$ by $\langle2\rangle1$ and
 $\sigma$-weak continuity of $\Ad\Gm(\Wt)$); von~Neumann density theorem for
 $\cM=\mathrm{CAR}(\mathfrak h_I)''$}
\lqed{2}{4}\lby{$\langle2\rangle1$--$\langle2\rangle3$; $\norm\Psi=\norm\Om=1$}
\lstep{1}{2}{$\mathrm F(\omega|_\cM,\omega^\beta|_\cM)\ge\abs{\ip\Om\Psi}=
 \abs{\ip\Om{\Gm(\Wt)\Om}}$.}
\lby{Alberti--Uhlmann 2002 (\texttt{math-ph/0202038}) Thm.~2(3), elementary half: the fidelity is
 the supremum of $\abs{\ip\phi\psi}$ over vector representatives, and $\Om$, $\Psi$ are vector
 representatives of $\omega|_\cM$, $\omega^\beta|_\cM$ by $\langle1\rangle1$ (transcribed with a
 screenshot in PA \S3, citedbox ``Alberti--Uhlmann 2002, Theorem 2(3)''); the last equality is
 $\ip\Om{T^*\Om}=\overline{\ip\Om{T\Om}}$}
\lstep{1}{3}{$\abs{\ip\Om{\Gm(\Wt)\Om}}=\det(\one-V^*V)^{1/2}$.}
\lby{PA Thm.~4.1 (vacuum overlap of a Bogoliubov unitary), whose hypotheses
 $V=\Pperp\Wt P\in\Stwo$, $A=P\Wt\Pperp\in\Stwo$, $\norm V<1$, $\norm A<1$ hold: $V\in\Stwo$ by
 Def.~\ref{def:adm}; $A$ and $V$ have the same singular values, because in the
 cosine--sine (CS) decomposition of the unitary $\Wt$ with respect to
 $H=PH\oplus\Pperp H$ the two off-diagonal blocks are the two ``sine'' blocks
 (\texttt{rate\_of\_remainder.tex} Thm.~2.3 $\langle1\rangle4$), so $A\in\Stwo$ and
 $\norm A=\norm V\le\norm V_2<1$}
\lstep{1}{4}{$\det(\one-V^*V)\ge\one-\norm V_2^2$, and
 $-\tfrac12\log\det(\one-V^*V)=\tfrac12\norm V_2^2+O(\norm V^4_2)$.}
\lby{\texttt{rate\_of\_remainder.tex} Thm.~2.5 $\langle1\rangle2$: $\det(\one-V^*V)=
 \prod_j(1-\sigma_j^2)\ge1-\sum_j\sigma_j^2=1-\norm V_2^2$; for the expansion,
 $-\tfrac12\log\prod_j(1-\sigma_j^2)=\tfrac12\sum_j\sigma_j^2+\tfrac14\sum_j\sigma_j^4+\ldots
 \in[\tfrac12\norm V_2^2,\tfrac12\norm V_2^2(1-\norm V_2^2)^{-1}]$}
\lqed{1}{5}\lby{$\langle1\rangle1$--$\langle1\rangle4$ and $-\log$ is decreasing}

\begin{verdictok}
Proposition~\ref{prop:uhl} is the brief's ``key point'' in precise form, and it is the only place
where the fidelity is touched: from here on everything is a statement about the single
Hilbert--Schmidt operator $V=\Pperp\Wt P$.  Two remarks.  (i) Unlike PA \S3 we do not need the
commutant unitary $u'=\Gm(\widehat W')^*$ explicitly: by Lemma~\ref{lem:torsor} the freedom
$u'$ provides is \emph{exactly} the freedom of choosing $\Wt$ inside $\mathcal E(W_I)$, and
Prop.~\ref{prop:uhl} is stated for each $\Wt$ separately.  (The two formulations agree:
$\Gm(\Wt_0\widehat W')=\Gm(\Wt_0)\Gm(\widehat W')$ up to a phase, and $\Gm(\widehat W')\in\cM'$
by PA Lemma~3.2.)  (ii) $\Wt$ does \emph{not} preserve $\mathfrak h_{I^c}$: since $W_I$ is an
isometry of $\mathfrak h_I$ \emph{onto the proper subspace} $\mathfrak h_{AB}$, an admissible
$\Wt$ maps $\mathfrak h_{I^c}$ onto $\mathfrak h_{AB}^\perp\supsetneq\mathfrak h_{I^c}$.  The
sentence ``hence $\Wt$ preserves $L^2(I^c)$'' in the Phase-4 brief is therefore false as stated,
and with it the idea of writing $\Wt=\exp(Z_I\oplus Z_e)$ with a block-diagonal generator; what
\emph{is} true and what is used below is Lemma~\ref{lem:torsor}: the admissible extensions form a
right torsor under the group of unitaries $\one_{\mathfrak h_I}\oplus W'$, whose generators
$\widehat Z_e=0\oplus Z_e$ \emph{are} block diagonal.  This is repaired once and for all in
\S\ref{sec:second}.
\end{verdictok}

\section{Working notes}\label{sec:notes}

\emph{(Reasoning recorded as it was produced; the polished statements are
\S\S\ref{sec:two}--\ref{sec:main}.)}

\textbf{(N1) The two orthogonality conditions.}  For $B\in\cH$ (so $B=\Pperp BP$) and
$\ell=\ell^*=E_I\ell E_I$, $\Tr(B^*\Pperp\ell P)=\Tr(B^*\ell)$, hence
$\ip B{\Pperp\ell P}_r=\tfrac12\Tr\bigl((B+B^*)\ell\bigr)=\tfrac12\Tr\bigl(E_I(B+B^*)E_I\,\ell\bigr)$;
for $Z_e^*=-Z_e=E_{I^c}Z_eE_{I^c}$, $\ip B{\Pperp Z_eP}_r=\tfrac12\Tr\bigl((B^*-B)Z_e\bigr)
=\tfrac12\Tr\bigl(E_{I^c}(B^*-B)E_{I^c}Z_e\bigr)$.  So $B\perp\Nvis$ iff $E_I(B+B^*)E_I=0$ and
$B\perp\Ext$ iff $E_{I^c}(B-B^*)E_{I^c}=0$.  With $\Xi:=B+B^*$ (self-adjoint, $P$-off-diagonal)
one has $B-B^*=(\one-2P)\Xi$, so \emph{(Q) is equivalent to}: a self-adjoint
$\Xi\in\Stwo$ with $P\Xi P=\Pperp\Xi\Pperp=0$, $E_I\Xi E_I=0$ and
$E_{I^c}(\one-2P)\Xi E_{I^c}=0$ vanishes.

\textbf{(N2) Halmos.}  In generic position write $H\cong\cK\oplus\cK$ with
$E_I=\mathrm{diag}(1,0)$ and $2P-\one=\begin{psmallmatrix}\gamma&\sigma\\ \sigma&-\gamma\end{psmallmatrix}$,
$\gamma=\cos2\Theta$, $\sigma=\sin2\Theta$.  Then $E_I\Xi E_I=0$ kills the $(1,1)$ entry, and
with $\Xi=\begin{psmallmatrix}0&w^*\\ w&2p\end{psmallmatrix}$ the three remaining conditions read
$\sigma w+w^*\sigma=0$, $[\gamma,w]=2p\sigma$, $\sigma w^*=2\gamma p$.  Multiplying the second by
$\sigma$ on the right and using $w\sigma=2p\gamma$ gives $\gamma p\gamma=p$, hence
$p=\gamma^np\gamma^n\to0$ because $\gamma^n\to0$ strongly ($\ker(\one-\gamma^2)=\ker\sigma^2=0$)
and $p$ is compact; then $\sigma w^*=0$ and $\sigma$ injective give $w=0$.  This is
\S\ref{sec:two}.  The only inputs are: generic position of $(P,E_I)$ (F.~and~M.~Riesz) and
$\Xi$ compact.  Compactness is essential: for $p$ merely bounded, $\gamma p\gamma=p$ has
solutions as soon as $1\in\operatorname{spec}\gamma$ (approximate eigenvectors), which is exactly
the ``continuous spectrum accumulating at the endpoints'' caveat; this is why the theorem is a
statement about $\Stwo$ and why $\overline\Ext=\Nvis^\perp$ can only be a \emph{closure}
statement.

\textbf{(N3) Second, independent route (duality).}  $\Nvis$ and $\Ext$ are, in the
one-particle--one-hole sector, the standard subspaces of $\cM=\cF(I)$ and of $\cM'$:
PA Lemma~7.5 (``Wick reduction'') proves $H_\cM\cap\cH^{(1,1)}=\overline{\Nvis}$ and
$H_{\cM'}\cap\cH^{(1,1)}=\overline{\mathcal K}$ with
$\mathcal K=\{\Pperp h'P:h'=h'^*\ \text{supported in }I^c\}$, i.e.\ $i\mathcal K=\Ext$.  For any
closed real subspace $K$ of a complex Hilbert space, $K^{\perp_{\mathbb R}}=iK'$ with $K'$ the
symplectic complement, and $K'=\overline{\cM'_{\rm sa}\Om}$ when $K=\overline{\cM_{\rm sa}\Om}$
(Tomita--Takesaki).  Compressing with the (complex-linear, orthogonal) projection $E$ onto
$\cH^{(1,1)}$ gives $\Nvis^\perp=i\overline{\mathcal K}=\overline\Ext$ again.  So (Q) is
\emph{equivalent to Haag (twisted) duality} for the free chiral fermion, read in the two-particle
sector.  Recorded as Theorem~\ref{thm:Qdual}.  The Halmos route is preferred as the primary proof
because it uses no net-theoretic input.

\textbf{(N4) Finite dimensions.}  $\dim_{\mathbb R}\Nvis=m^2$ ($m=\dim\mathfrak h_I$),
$\dim_{\mathbb R}\cH=2k(n-k)$, so $\dim\Nvis^\perp=2k(n-k)-m^2$, while the exterior
anti-Hermitian matrices form a space of dimension $(n-m)^2$; at half filling $k=n/2$ these agree
iff $m=n/2$, which is exactly the condition for $(P,E_I)$ to be in generic position in finite
dimensions.  E0's computation (card \texttt{finite-dim-exterior-density}) finds $\Ext=\Nvis^\perp$
also for $m\ne n/2$, where $Z_e\mapsto\Pperp Z_eP$ then has a kernel; this is consistent but is
\emph{not} what the continuum proof uses.

\section{The visible space, the exterior space, and the two lemmas}\label{sec:lemmas}

\begin{definition}[Visible and exterior directions]\label{def:NE}
Inside the real Hilbert space $\cH=\Stwo(PH,\Pperp H)$ of \eqref{eq:realHS} put
\begin{align}
 \Nvis&:=\bigl\{\Pperp\ell P:\ \ell=\ell^*\ \text{finite rank},\ \ell=E_I\ell E_I\bigr\},
 \label{eq:N}\\
 \Ext&:=\bigl\{\Pperp Z_eP:\ Z_e^*=-Z_e\ \text{bounded},\ Z_e=E_{I^c}Z_eE_{I^c},\
 \Pperp Z_eP\in\Stwo\bigr\},\label{eq:E}\\
 \Ext_{\rm fr}&:=\bigl\{\Pperp Z_eP:\ Z_e^*=-Z_e\ \text{finite rank},\ Z_e=E_{I^c}Z_eE_{I^c}\bigr\}
 \ \subseteq\ \Ext .\label{eq:Efr}
\end{align}
$\Pi$ is the $\ip\cdot\cdot_r$-orthogonal projection of $\cH$ onto $\overline{\Nvis}$.
\end{definition}

\begin{lemma}[L1: the pairing identity, and invisibility of exterior directions]\label{lem:L1}
Let $Z$ be bounded and anti-self-adjoint with $\Pperp ZP\in\Stwo$, and write $M_Z:=\Pperp ZP$.
\begin{enumerate}[label=(\alph*),leftmargin=2.0em,itemsep=1pt]
\item For every finite-rank $\ell=\ell^*=E_I\ell E_I$,
 \begin{equation}\label{eq:L1}
  \bigl\langle \Pperp\ell P,\ M_Z\bigr\rangle_r=-\tfrac12\Tr\bigl(\ell\,E_I[P,Z]E_I\bigr).
 \end{equation}
\item $M_Z\perp\Nvis$ if and only if $E_I[P,Z]E_I=0$.
\item If $Z=Z_e$ is exterior ($Z_e=E_{I^c}Z_eE_{I^c}$) then $E_I[P,Z_e]E_I=0$; hence
 $\overline{\Ext}\subseteq\Nvis^{\perp}$.
\item For $B\in\cH$: $B\perp\Nvis\iff E_I(B+B^*)E_I=0$, and
 $B\perp\Ext\iff E_{I^c}(B-B^*)E_{I^c}=0$.
\end{enumerate}
\end{lemma}
\lstep{1}{1}{$\ip{\Pperp\ell P}{M_Z}_r=\Re\Tr\bigl(\ell\,\Pperp ZP\bigr)$.}
\lby{$\ip XY_r=\Re\Tr(X^*Y)$ with $X=\Pperp\ell P$, $X^*=P\ell\Pperp$, so
 $\Tr(X^*Y)=\Tr(P\ell\Pperp ZP)=\Tr(\ell\Pperp ZP)$ by cyclicity and $P^2=P$; the trace exists
 because $\ell$ is finite rank}
\lstep{1}{2}{$\overline{\Tr(\ell\Pperp ZP)}=-\Tr(\ell\,PZ\Pperp)$, hence
 $\ip{\Pperp\ell P}{M_Z}_r=\tfrac12\Tr\bigl(\ell\,(\Pperp ZP-PZ\Pperp)\bigr)$.}
\lby{$\overline{\Tr T}=\Tr(T^*)$, $(\ell\Pperp ZP)^*=PZ^*\Pperp\ell=-PZ\Pperp\ell$ ($Z^*=-Z$,
 $\ell^*=\ell$), and cyclicity; then $\Re w=\tfrac12(w+\bar w)$ and $\langle1\rangle1$}
\lstep{1}{3}{(a).  $[P,Z]=PZ\Pperp-\Pperp ZP$, so $\Pperp ZP-PZ\Pperp=-[P,Z]$; and
 $\Tr(\ell[P,Z])=\Tr(\ell E_I[P,Z]E_I)$.}
\lby{$[P,Z]=PZ(P+\Pperp)-(P+\Pperp)ZP=PZ\Pperp-\Pperp ZP$; the second claim is
 $\ell=E_I\ell E_I$ and cyclicity; insert into $\langle1\rangle2$}
\lstep{1}{4}{(b).}
\lby{``$\Leftarrow$'' is immediate from (a).  ``$\Rightarrow$'': if the left side of \eqref{eq:L1}
 vanishes for all such $\ell$, take $\ell=\ket f\bra f$, $f\in\mathfrak h_I$, to get
 $\ip f{E_I[P,Z]E_If}=0$ for all $f$; and $E_I[P,Z]E_I$ is \emph{self}-adjoint, because
 $[P,Z]^*=Z^*P-PZ^*=-ZP+PZ=[P,Z]$ for $Z^*=-Z$; a self-adjoint operator all of whose diagonal
 matrix elements vanish is $0$}
\lstep{1}{5}{(c).}
\lby{$Z_e=E_{I^c}Z_eE_{I^c}$ and $E_IE_{I^c}=0$ give $E_IZ_e=Z_eE_I=0$, hence
 $E_I[P,Z_e]E_I=E_IPZ_eE_I-E_IZ_ePE_I=0$; then (b) gives $\Ext\perp\Nvis$, and an orthogonal
 complement is closed}
\lstep{1}{6}{(d).}
\lproof
\lstep{2}{1}{$\ip{\Pperp\ell P}B_r=\tfrac12\Tr\bigl(E_I(B+B^*)E_I\,\ell\bigr)$ and
 $\ip{\Pperp Z_eP}B_r=-\tfrac12\Tr\bigl(Z_e\,E_{I^c}(B-B^*)E_{I^c}\bigr)$.}
\lby{First: $(\Pperp\ell P)^*B=P\ell\Pperp B$ and $B=\Pperp BP$ give
 $\Tr((\Pperp\ell P)^*B)=\Tr(\ell B)$; $\overline{\Tr(\ell B)}=\Tr(B^*\ell^*)=\Tr(\ell B^*)$, so
 the real part is $\tfrac12\Tr(\ell(B+B^*))=\tfrac12\Tr(E_I(B+B^*)E_I\ell)$ by $\ell=E_I\ell E_I$.
 Second: $(\Pperp Z_eP)^*B=PZ_e^*\Pperp B$ gives $\Tr((\Pperp Z_eP)^*B)=\Tr(Z_e^*B)=-\Tr(Z_eB)$;
 $\overline{\Tr(Z_eB)}=\Tr(B^*Z_e^*)=-\Tr(Z_eB^*)$, so
 $\Re\Tr(Z_eB)=\tfrac12\Tr(Z_e(B-B^*))$, and $Z_e=E_{I^c}Z_eE_{I^c}$ inserts the two $E_{I^c}$}
\lstep{2}{2}{The vanishing of the first expression for all finite-rank self-adjoint
 $\ell$ on $\mathfrak h_I$ is equivalent to $E_I(B+B^*)E_I=0$; the vanishing of the second for all
 finite-rank anti-self-adjoint $Z_e$ on $\mathfrak h_{I^c}$ is equivalent to
 $E_{I^c}(B-B^*)E_{I^c}=0$.}
\lby{$E_I(B+B^*)E_I$ is self-adjoint and Hilbert--Schmidt: test with $\ell=\ket f\bra f$ and use
 $\langle1\rangle4$'s polarisation argument.  $A:=E_{I^c}(B-B^*)E_{I^c}$ is anti-self-adjoint and
 Hilbert--Schmidt; the finite-rank anti-self-adjoint operators on $\mathfrak h_{I^c}$ are
 $\Stwo$-dense among the anti-self-adjoint elements of $\Stwo(\mathfrak h_{I^c})$ and
 $Z_e\mapsto\Tr(Z_eA)$ is $\Stwo$-continuous, so the hypothesis extends to $Z_e=-A$, giving
 $-\Tr(AA)=\Tr(A^*A)=\norm A_2^2=0$}
\lqed{2}{3}\lby{$\langle2\rangle1$, $\langle2\rangle2$}
\lqed{1}{7}\lby{$\langle1\rangle3$--$\langle1\rangle6$}

\begin{lemma}[L2: the Legendre form of a commutator defect]\label{lem:L2}
Let $Z$ be bounded and anti-self-adjoint with $M_Z=\Pperp ZP\in\Stwo$, and put
$\delta_Z:=E_I[P,Z]E_I$ (a self-adjoint Hilbert--Schmidt operator on $\mathfrak h_I$).  Then
\begin{equation}\label{eq:L2}
 g_Q(\delta_Z)=\tfrac12\norm{\Pi M_Z}_2^2 .
\end{equation}
\end{lemma}
\lstep{1}{1}{For self-adjoint finite-rank $t=E_ItE_I$, $\Tr\bigl(tQt(\one-Q)\bigr)=\norm{N}_2^2$
 with $N:=\Pperp tP\in\Nvis$, and $t\mapsto N$ maps the set of such $t$ \emph{onto} $\Nvis$.}
\lby{S10 Lemma~4.1 $\langle1\rangle1$ (with the vacuum symbol $Q=E_IPE_I$, S10 eq.~(1)):
 $\Tr(tQt(\one-Q))=\Tr(PtP^\perp t)=\Tr(N^*N)$; surjectivity is Definition~\ref{def:NE}}
\lstep{1}{2}{$\Tr(t\delta_Z)=-2\ip N{M_Z}_r$.}
\lby{$\Tr(t\delta_Z)=\Tr(t[P,Z])$ since $t=E_ItE_I$; then Lemma~\ref{lem:L1}(a) with $\ell=t$}
\lstep{1}{3}{$8g_Q(\delta_Z)=\sup_{t}\dfrac{(\Tr t\delta_Z)^2}{\Tr(tQt(\one-Q))}
 =4\sup_{0\ne N\in\Nvis}\dfrac{\ip N{M_Z}_r^2}{\norm N_2^2}=4\norm{\Pi M_Z}_2^2$.}
\lby{\eqref{eq:defect} for the first equality (the supremum over self-adjoint $t$ may be
 restricted to finite-rank $t$: both $t\mapsto\Tr(t\delta_Z)$ and $t\mapsto\Tr(tQt(1-Q))$ are
 continuous in $\norm\cdot_2$ and finite-rank operators are dense, and the quotient is
 scale-invariant); $\langle1\rangle1$, $\langle1\rangle2$ for the second; and for the third the
 elementary Hilbert-space identity $\sup_{0\ne N\in\Nvis}\ip NM_r^2/\norm N^2=\norm{\Pi M}^2$
 (S10 Prop.~4.3 $\langle1\rangle2$), valid in the real Hilbert space $(\cH,\ip\cdot\cdot_r)$ with
 $\Pi$ the orthogonal projection onto $\overline{\Nvis}$}
\lqed{1}{4}\lby{$\langle1\rangle3$}

\begin{verdictok}
Lemmas~\ref{lem:L1} and \ref{lem:L2} are exactly (L1) and (L2) of the Phase-4 brief and are
proved here in the generality needed: $Z$ is an arbitrary bounded anti-self-adjoint operator with
$\Pperp ZP\in\Stwo$, not necessarily a vector field.  (For unbounded generators such as
$Z=-D_\chi$ the two lemmas hold verbatim with $\ell$, $t$ finite rank, because only
$\Tr(\ell[P,Z])$ and $\Pperp ZP\in\Stwo$ are used; boundedness enters nowhere in
$\langle1\rangle1$--$\langle1\rangle3$ of Lemma~\ref{lem:L1}.)  Lemma~\ref{lem:L2} is S10
Lemma~4.1/Prop.~4.3 with $sD_\chi$ replaced by $Z$; the only additional point checked here is
that restricting the supremum defining $g_Q$ to finite-rank $t$ does not change it.
\end{verdictok}

\section{The reduction theorem}\label{sec:reduction}

The generator of an admissible extension is \emph{not} block diagonal (verdict after
Prop.~\ref{prop:uhl}).  The correct linearised bookkeeping is the following.

\begin{lemma}[The prescribed column and the free block]\label{lem:column}
Let $\tau\mapsto\Wt_\tau=e^{\tau Z}$ be a one-parameter family of unitaries of $H$ with
$Z^*=-Z$, and suppose the restrictions $\Wt_\tau|_{\mathfrak h_I}=W_I(\tau)$ are prescribed.  Then
\begin{enumerate}[label=(\alph*),leftmargin=2.0em,itemsep=1pt]
\item the whole \emph{column} $ZE_I$ --- i.e.\ both blocks $E_IZE_I$ and $E_{I^c}ZE_I$ --- is
 determined by the channel data, namely $ZE_I=\frac{\dd}{\dd\tau}W_I(\tau)|_{\tau=0}$;
\item so is $E_IZE_{I^c}=-(E_{I^c}ZE_I)^*$;
\item at first order the only free block is $\widehat Z_e:=E_{I^c}ZE_{I^c}$: the map
 $\Wt_0\mapsto\Wt_0e^{\tau Z_e}$ ($Z_e$ exterior, $\tau\in\mathbb R$) leaves every channel datum
 unchanged (Lemma~\ref{lem:torsor}) and shifts the generator by exactly $0\oplus Z_e$, and
 $\Pperp(0\oplus Z_e)P$ ranges over $\Ext$;
\item $E_I[P,Z]E_I$ is determined by the column alone, hence does not depend on $\widehat Z_e$.
\end{enumerate}
\end{lemma}
\lstep{1}{1}{(a), (b), (c).}
\lby{(a): differentiate $\Wt_\tau E_I=W_I(\tau)E_I$ at $\tau=0$.  (b): $Z^*=-Z$ gives
 $E_IZE_{I^c}=-(E_{I^c}Z^*E_I)^*\cdot(-1)^{0}$, explicitly
 $(E_{I^c}ZE_I)^*=E_IZ^*E_{I^c}=-E_IZE_{I^c}$.  (c): the four blocks exhaust $Z$, the first three
 are fixed by (a),(b), and $\widehat Z_e$ is anti-self-adjoint and otherwise unconstrained; the
 last clause is Definition~\ref{def:NE}}
\lstep{1}{2}{(d).}
\lby{$E_I[P,Z]E_I=E_IPZE_I-E_IZPE_I$; the first term uses only the column $ZE_I$, and
 $E_IZPE_I=(E_IZE_I)PE_I+(E_IZE_{I^c})PE_I$ uses only blocks fixed by (a),(b)}
\lqed{1}{3}\lby{$\langle1\rangle1$, $\langle1\rangle2$}

\begin{remark}[Do not truncate]\label{rem:trunc}
Lemma~\ref{lem:column} is the structural correction recorded by track~E2 (card
\texttt{cft\_cmi:e2-pre\-scribed-block-is-the-column}): for the zero-collar compression the field
$\chi$ does not vanish at the outer endpoint of $I$, the flow pushes points across $\partial I$,
and $W_I$ is an isometry of $\mathfrak h_I$ that is \emph{not} onto.  Consequently ``$Z_I$'' in the
Phase-4 brief must be read as the \emph{column}, e.g.\ $Z_0=-D_\chi$ with the full line vector
field, and \emph{not} as the truncation $E_ID_\chi E_I$: E2 measures
$\norm{\Pi\Pperp(E_ID_\chi E_I)P}_2^2/\norm u_2^2=2608,6188,11052,25290$ at $L=128,192,256,384$
(divergent like $L^2$: sharp truncation of a derivative creates an $O(L)$ boundary term), so the
truncation is not the generator of any admissible channel.  Theorem~\ref{thm:reduction} below is
stated for an arbitrary $Z_0$ with $\Pperp Z_0P\in\Stwo$ and is applied with $Z_0$ a genuine
column; by Corollary~\ref{cor:extindep} the value does not depend on which one.
\end{remark}

\begin{theorem}[Optimising the Uhlmann bound over exterior extensions]\label{thm:reduction}
Let $Z_0$ be anti-self-adjoint (possibly unbounded) with $M_0:=\Pperp Z_0P\in\Stwo$, and
$\delta_{Z_0}=E_I[P,Z_0]E_I$.  Then
\begin{equation}\label{eq:reduction}
 \inf_{Z_e}\ \tfrac12\bigl\|\Pperp(Z_0+Z_e)P\bigr\|_2^2
 \ =\ \tfrac12\dist\bigl(M_0,\Ext\bigr)^2
 \ =\ \tfrac12\norm{\Pi M_0}_2^2+\tfrac12\dist\bigl((\one-\Pi)M_0,\Ext\bigr)^2
 \ \ge\ g_Q(\delta_{Z_0}),
\end{equation}
the infimum being over all exterior anti-self-adjoint $Z_e$ as in \eqref{eq:E}, with
\begin{equation}\label{eq:iff}
 \inf_{Z_e}\ \tfrac12\bigl\|\Pperp(Z_0+Z_e)P\bigr\|_2^2=g_Q(\delta_{Z_0})
 \qquad\Longleftrightarrow\qquad
 (\one-\Pi)M_0\in\overline{\Ext}.
\end{equation}
\end{theorem}
\lstep{1}{1}{$\inf_{Z_e}\norm{M_0+\Pperp Z_eP}_2^2=\dist(M_0,\Ext)^2=\dist(M_0,\overline\Ext)^2$.}
\lby{$\Ext$ is a real-linear subspace of $\cH$ (Definition~\ref{def:NE}; $Z_e\mapsto\Pperp Z_eP$
 is real-linear and the exterior anti-self-adjoint operators form a real vector space), so
 $\{M_0+\Pperp Z_eP\}=M_0+\Ext$; and the distance to a set equals the distance to its closure}
\lstep{1}{2}{$\dist(M_0,\overline\Ext)^2=\norm{\Pi M_0}_2^2+\dist\bigl((\one-\Pi)M_0,\overline\Ext\bigr)^2$.}
\lby{$\overline\Ext\subseteq\Nvis^\perp=\ran(\one-\Pi)$ by Lemma~\ref{lem:L1}(c), so for
 $Y\in\overline\Ext$ the Pythagoras identity
 $\norm{M_0-Y}^2=\norm{\Pi M_0}^2+\norm{(\one-\Pi)M_0-Y}^2$ holds; take the infimum over $Y$}
\lstep{1}{3}{$\tfrac12\norm{\Pi M_0}_2^2=g_Q(\delta_{Z_0})$.}
\lby{Lemma~\ref{lem:L2} (its proof used boundedness of $Z$ nowhere, see the verdict above)}
\lstep{1}{4}{\eqref{eq:reduction} and \eqref{eq:iff}.}
\lby{$\langle1\rangle1$--$\langle1\rangle3$; the infimum equals $g_Q$ iff the second term of
 $\langle1\rangle2$ vanishes, i.e.\ iff $(\one-\Pi)M_0\in\overline\Ext$ (a closed set)}
\lqed{1}{5}\lby{$\langle1\rangle4$}

\begin{corollary}[Extension independence]\label{cor:extindep}
If $Z_0,Z_0'$ are anti-self-adjoint with $\Pperp Z_0P,\Pperp Z_0'P\in\Stwo$ and
$E_I[P,Z_0]E_I=E_I[P,Z_0']E_I$, then $\dist(\Pperp Z_0P,\Ext)=\dist(\Pperp Z_0'P,\Ext)$ and
$\Pi\Pperp Z_0P=\Pi\Pperp Z_0'P$: the optimised Uhlmann bound depends on the channel only through
its first-order defect.
\end{corollary}
\lby{$Z_0-Z_0'$ is anti-self-adjoint with $E_I[P,Z_0-Z_0']E_I=0$, so
$\Pperp(Z_0-Z_0')P\in\Nvis^\perp$ by Lemma~\ref{lem:L1}(b); by Theorem~\ref{thm:Q}
$\Nvis^\perp=\overline\Ext$, so the two elements differ by an element of the closed subspace
$\overline\Ext$, which changes neither the distance to $\overline\Ext$ nor the component
$\Pi(\cdot)$ (as $\overline\Ext\perp\overline\Nvis$)}

\begin{corollary}[What (UH) needs]\label{cor:UHneeds}
For the family of Lemma~6.1 of S10, $Z_0=-\bigl(sD_\chi+\eps\Gs_*\bigr)$ up to
$O(s^2)+O(\eps^2)$, the leading-order Uhlmann bound optimised over exterior extensions equals
$g_Q(\delta(\eps))$ if and only if $(\one-\Pi)\Pperp(sD_\chi+\eps\Gs_*)P\in\overline\Ext$.  A
sufficient condition, independent of the channel, is $\overline\Ext=\Nvis^{\perp}$, which is
question (Q).
\end{corollary}
\lby{Theorem~\ref{thm:reduction} applied to $Z_0$, S10 Lemma~6.1(2)--(3) for the identification of
 $\delta(\eps)$ and of $g_Q(\delta(\eps))$, and Lemma~\ref{lem:L1}(c) for
 $\overline\Ext\subseteq\Nvis^\perp$}

\section{Question (Q): the exterior directions exhaust the invisible space}\label{sec:two}

\begin{lemma}[Generic position of $(P,E_I)$]\label{lem:generic}
$\mathfrak h_I\cap PH=\mathfrak h_I\cap\Pperp H=\mathfrak h_{I^c}\cap PH
=\mathfrak h_{I^c}\cap\Pperp H=\{0\}$.  Consequently
\begin{enumerate}[label=(\alph*),leftmargin=2.0em,itemsep=1pt]
\item $\Gamma':=E_{I^c}(2P-\one)E_{I^c}$, viewed as a self-adjoint contraction of
 $\mathfrak h_{I^c}$, has no eigenvalue $+1$ and no eigenvalue $-1$; hence
 $\Gamma'^{\,n}\to0$ strongly;
\item $E_{I^c}PE_I:\mathfrak h_I\to\mathfrak h_{I^c}$ is injective.
\end{enumerate}
\end{lemma}
\lstep{1}{1}{$PH=\{f\in L^2(\mathbb R):\hat f=0\ \text{a.e.\ on a half-line}\}$ is the boundary-value
 Hardy space $H^2$ of a half-plane, and a nonzero $f\in H^2$ cannot vanish on a set of positive
 Lebesgue measure.}
\lby{Paley--Wiener (W.~Rudin, \emph{Real and Complex Analysis}, 3rd ed., Thm.~19.2, p.~372):
 $L^2$ functions with Fourier transform supported in a half-line are exactly the boundary values of
 $H^2$ functions of the corresponding half-plane; and Rudin, Thm.~17.18, p.~344 ($\log\abs{f^*}\in
 L^1(\mathbb R)$ for $0\ne f\in H^p$), which forces $f^*\ne0$ a.e.  \emph{Screenshots NOT OBTAINED}:
 these are textbook results and the book is not in \texttt{refs/}; both statements are classical
 (F.~and~M.~Riesz / Szeg\H o) and are quoted with page-exact locators}
\lstep{1}{2}{The four intersections vanish.}
\lby{$I$ and $I^c$ both have positive Lebesgue measure (both are nonempty open sets); a nonzero
 $f\in PH$ vanishing a.e.\ on $I$ or on $I^c$ contradicts $\langle1\rangle1$; for $\Pperp H$ use
 $\overline{\Pperp H}=\Pperp H$ is the Hardy space of the \emph{other} half-plane and the same
 argument (or apply $\langle1\rangle1$ to $\bar f$)}
\lstep{1}{3}{(a).}
\lproof
\lstep{2}{1}{If $\Gamma'\xi=\xi$ with $\xi\in\mathfrak h_{I^c}$, $\norm\xi=1$, then $\norm{P\xi}=1$
 and hence $\xi\in\mathfrak h_{I^c}\cap PH=\{0\}$, a contradiction; the case $\Gamma'\xi=-\xi$ is
 the same with $\Pperp$.}
\lby{$1=\ip\xi{\Gamma'\xi}=\ip\xi{(2P-\one)\xi}=2\norm{P\xi}^2-1$ for $\xi\in\mathfrak h_{I^c}$,
 so $\norm{P\xi}^2=1=\norm\xi^2$ and $\Pperp\xi=0$; then $\langle1\rangle2$}
\lstep{2}{2}{$\Gamma'^{\,n}\to0$ strongly.}
\lby{$\norm{\Gamma'^{\,n}\xi}^2=\int_{[-1,1]}\lambda^{2n}\dd\mu_\xi(\lambda)\to\mu_\xi(\{-1,1\})
 =\norm{\mathbf 1_{\{-1,1\}}(\Gamma')\xi}^2=0$ by dominated convergence ($\abs\lambda\le1$) and
 $\langle2\rangle1$, the spectral projection of a self-adjoint operator at a point being the
 eigenprojection at that point}
\lqed{2}{3}\lby{$\langle2\rangle1$, $\langle2\rangle2$}
\lstep{1}{4}{(b).}
\lby{for $\xi\in\mathfrak h_I$, $\norm{E_{I^c}PE_I\xi}^2=\ip\xi{E_IPE_{I^c}PE_I\xi}
 =\ip\xi{(A-A^2)\xi}$ with $A:=E_IPE_I|_{\mathfrak h_I}$ (insert $E_{I^c}=\one-E_I$ and use
 $P^2=P$); $A$ and $\one-A$ are commuting positive contractions, so $A-A^2=0$ on $\xi$ iff
 $A^{1/2}(\one-A)^{1/2}\xi=0$ iff $\xi\in\ker A+\ker(\one-A)$; and $\ker A=\mathfrak h_I\cap\Pperp H
 =\{0\}$ (as $\ip\xi{A\xi}=\norm{P\xi}^2$), $\ker(\one-A)=\mathfrak h_I\cap PH=\{0\}$, by
 $\langle1\rangle2$}
\lqed{1}{5}\lby{$\langle1\rangle2$--$\langle1\rangle4$}

\begin{remark}[Cross-check of Lemma~\ref{lem:generic}]\label{rem:BJL}
Generic position of the pair (one-particle vacuum projection, localisation projection) is exactly
the hypothesis under which Baumg\"artel--Jurke--Lled\'o 2002 (\texttt{math-ph/0204029}) run their
modular theory, and by their Prop.~3.4 it is \emph{equivalent} to $\Om$ being cyclic and separating
for the local algebra, i.e.\ to Reeh--Schlieder (screenshot \texttt{shots/BJL\_p7.png}, p.~10,
already used in \texttt{optimality\_all\_channels.tex} \S3 $\langle1\rangle1$).  Lemma~\ref{lem:generic}
is thus the one-particle form of Reeh--Schlieder for $\cF(I)$, proved here directly.
\end{remark}

\begin{theorem}[(Q) is true: $\overline{\Ext}=\Nvis^{\perp}$]\label{thm:Q}
In the real Hilbert space $\cH=\Stwo(PH,\Pperp H)$ of \eqref{eq:realHS},
\begin{equation}\label{eq:Q}
 \overline{\Ext_{\rm fr}}=\overline{\Ext}=\Nvis^{\perp}
 =\bigl\{B\in\cH:\ E_I(B+B^*)E_I=0\bigr\},
\end{equation}
and dually $\overline{\Nvis}=\Ext^{\perp}$.  Equivalently: every direction of $\cH$ that is
invisible to the interior $I$ is an $\cH$-limit of exterior directions $\Pperp Z_eP$ with
$Z_e^*=-Z_e$ supported in $I^c$, and the exterior ones can be taken finite rank.
\end{theorem}
\lstep{1}{1}{\lsuffices \lprove: if $B\in\cH$ satisfies $B\perp\Nvis$ and $B\perp\Ext_{\rm fr}$
 then $B=0$.}
\lby{$\Ext_{\rm fr}\subseteq\Ext\subseteq\Nvis^\perp$ (Definition~\ref{def:NE} and
 Lemma~\ref{lem:L1}(c)), so all three closures sit inside the closed real subspace $\Nvis^\perp$;
 the assertion says that the orthogonal complement of $\overline{\Ext_{\rm fr}}$ \emph{inside}
 $\Nvis^\perp$ is $\{0\}$, which for a closed subspace of a Hilbert space gives
 $\overline{\Ext_{\rm fr}}=\Nvis^\perp$, and then the two intermediate closures are squeezed;
 the last description of $\Nvis^\perp$ is Lemma~\ref{lem:L1}(d), and $\overline\Nvis=\Ext^\perp$
 follows by taking orthogonal complements}
\lstep{1}{2}{\lassume $B\in\cH$, $B\perp\Nvis$, $B\perp\Ext_{\rm fr}$.
 \ldefine $\Xi:=B+B^*$, $\Gamma:=2P-\one$, $E:=E_I$, $E^c:=E_{I^c}$,
 $\Gamma':=E^c\Gamma E^c$, $\pi:=E^c\Xi E^c$.  Then $\Xi=\Xi^*\in\Stwo$ and
 \begin{equation}\label{eq:threeconds}
  \text{(i) }\Gamma\Xi+\Xi\Gamma=0,\qquad\text{(ii) }E\Xi E=0,\qquad\text{(iii) }E^c\Gamma\Xi E^c=0 .
 \end{equation}}
\lby{(i): $B=\Pperp BP$ gives $P\Xi P=\Pperp\Xi\Pperp=0$, i.e.\ $\Gamma\Xi\Gamma=-\Xi$
 ($\Gamma$ acts as $+1$ on $PH$ and as $-1$ on $\Pperp H$, so conjugation by $\Gamma$ changes the
 sign of the $P$-off-diagonal part), and multiplying by $\Gamma$ on the right gives (i).
 (ii): Lemma~\ref{lem:L1}(d), first item.  (iii): Lemma~\ref{lem:L1}(d), second item, gives
 $E^c(B-B^*)E^c=0$ (the finite-rank exterior operators suffice: see its proof
 $\langle2\rangle2$), and $B-B^*=-\Gamma\Xi$ because $\Gamma P=P$, $\Gamma\Pperp=-\Pperp$ give
 $\Gamma\Xi=\Gamma\Pperp BP+\Gamma PB^*\Pperp=-\Pperp BP+PB^*\Pperp=-(B-B^*)$}

\lstep{1}{3}{$\Gamma'\pi\Gamma'=\pi$.}
\lproof
\lstep{2}{1}{$E^c\Gamma\Xi\Gamma E^c=-\pi$.}
\lby{(i) gives $\Gamma\Xi\Gamma=-\Xi\Gamma^2=-\Xi$; compress with $E^c$}
\lstep{2}{2}{$E^c\Gamma E\,\Xi E^c=-\Gamma'\pi$ and $E^c\Xi\,E\Gamma E^c=-\pi\Gamma'$.}
\lby{insert $\one=E+E^c$ between $\Gamma$ and $\Xi$ in (iii):
 $E^c\Gamma E\Xi E^c+E^c\Gamma E^c\Xi E^c=0$, and $E^c\Gamma E^c\Xi E^c=\Gamma'\pi$ because
 $E^cE^c=E^c$; the second identity is the adjoint of the first, $\Xi$ and $\Gamma$ being
 self-adjoint}
\lstep{2}{3}{$E^c\Gamma\Xi\Gamma E^c=0-\Gamma'\pi\Gamma'-\Gamma'\pi\Gamma'+\Gamma'\pi\Gamma'
 =-\Gamma'\pi\Gamma'$.}
\lby{insert $\one=E+E^c$ twice, obtaining the four terms
 $E^c\Gamma E\Xi E\Gamma E^c$, $E^c\Gamma E\Xi E^c\Gamma E^c$, $E^c\Gamma E^c\Xi E\Gamma E^c$,
 $E^c\Gamma E^c\Xi E^c\Gamma E^c$; the first vanishes by (ii), the second is
 $(E^c\Gamma E\Xi E^c)\Gamma'=-\Gamma'\pi\Gamma'$ and the third is
 $\Gamma'(E^c\Xi E\Gamma E^c)=-\Gamma'\pi\Gamma'$ by $\langle2\rangle2$, and the fourth is
 $\Gamma'\pi\Gamma'$}
\lqed{2}{4}\lby{$\langle2\rangle1$ and $\langle2\rangle3$}
\lstep{1}{4}{$\pi=0$.}
\lby{iterating $\langle1\rangle3$ gives $\pi=\Gamma'^{\,n}\pi\Gamma'^{\,n}$ for all $n\ge1$;
 $\pi$ is compact (it is Hilbert--Schmidt, being a compression of $\Xi=B+B^*\in\Stwo$) and
 $\Gamma'^{\,n}\to0$ strongly with $\norm{\Gamma'}\le1$ (Lemma~\ref{lem:generic}(a)), whence
 $\norm{\Gamma'^{\,n}\pi}\to0$ in operator norm (for compact $\pi$, strong convergence is uniform
 on the precompact set $\pi(\text{unit ball})$), and therefore
 $\norm\pi\le\norm{\Gamma'^{\,n}\pi}\,\norm{\Gamma'^{\,n}}\to0$}
\lstep{1}{5}{$E\Xi E^c=0$ and $E^c\Xi E=0$.}
\lby{$\langle2\rangle2$ with $\pi=0$ gives $E^c\Gamma E\,\Xi E^c=0$; and
 $E^c\Gamma E=E^c(2P-\one)E=2E^cPE$ because $E^cE=0$, so
 $E^cPE\,(E\Xi E^c)=0$; by Lemma~\ref{lem:generic}(b) $E^cPE$ is injective, so
 $\ran(E\Xi E^c)\subseteq\ker(E^cPE)=\{0\}$; the second identity is the adjoint}
\lstep{1}{6}{$\Xi=0$, hence $B=0$.}
\lby{$\Xi=E\Xi E+E\Xi E^c+E^c\Xi E+E^c\Xi E^c=0+0+0+0$ by (ii), $\langle1\rangle5$,
 $\langle1\rangle4$; and $B=\Pperp\Xi P$ because $B=\Pperp BP$ and $B^*=PB^*\Pperp$}
\lqed{1}{7}\lby{$\langle1\rangle1$, $\langle1\rangle6$}

\begin{verdictok}
Theorem~\ref{thm:Q} answers question (Q) of the Phase-4 brief affirmatively, and the proof is
elementary and self-contained: the only input is generic position of $(P,E_I)$
(Lemma~\ref{lem:generic}, i.e.\ Hardy-space uniqueness, equivalently Reeh--Schlieder).
\emph{Compactness of $B$ is not needed} (referee report \S3): if $\pi$ is merely bounded and
$\pi=\Gamma'^{\,n}\pi\Gamma'^{\,n}$, then for all $\xi,\eta$
\[
 \ip\xi{\pi\eta}=\ip{\Gamma'^{\,n}\xi}{\pi\,\Gamma'^{\,n}\eta},\qquad
 \abs{\ip\xi{\pi\eta}}\le\norm{\Gamma'^{\,n}\xi}\,\norm\pi\,\norm{\Gamma'^{\,n}\eta}\to0,
\]
so $\pi=0$ by two-sided strong convergence alone; step $\langle1\rangle4$ may be read with
``compact'' deleted.  (The earlier version of this box claimed the opposite --- that
$\Gamma'\pi\Gamma'=\pi$ has nonzero bounded solutions built from approximate eigenvectors.  That
is false: approximate eigenvectors at $\pm1$ converge weakly to $0$, since no eigenvector
exists, so every weak-$*$ cluster point of $\ket{\xi_n}\bra{\xi_n}$ vanishes.  See
\S\ref{sec:corr} R3.)  Why then only a \emph{closure}?  Not because of the fixed-point
equation, but because $Z_e\mapsto\Pperp Z_eP$ has an unbounded inverse on $\Nvis^\perp$: E2
\S2.3 measures gains $\sqrt{d^c_{ij}/2}$ whose lower quartile falls from $10^{-2.5}$ at $L=96$
to $10^{-7.5}$ at $L=384$, i.e.\ a growing fraction of $\Nvis^\perp$ is reached only with
exponentially large $\norm{Z_e}$.  Generic position, by contrast, is used essentially:
  (a): the four
intersections enter in two groups --- $\mathfrak h_{I^c}\cap P^{\pm}H$ through
Lemma~\ref{lem:generic}(a) (used for $\pi=0$) and $\mathfrak h_I\cap P^{\pm}H$ through
Lemma~\ref{lem:generic}(b) (used for the off-diagonal blocks), $P^+:=P$, $P^-:=\Pperp$ --- and one
failure in \emph{each} group breaks the theorem: if $0\ne f\in\mathfrak h_{I^c}\cap PH$ and
$0\ne g\in\mathfrak h_I\cap\Pperp H$ then $B:=\ket g\bra f$ is a nonzero element of
$\Nvis^\perp\cap\Ext^\perp$ (indeed $B=\Pperp BP$, $E_If=0$ kills $E_I(B+B^*)E_I$, and
$E_{I^c}g=0$ kills $E_{I^c}(B-B^*)E_{I^c}$).  A failure in one group alone need not break it, and
does not in E2's model, where $Q$ has $z$ exact eigenvalues $0$ and $z$ exact eigenvalues $1$
(so both members of the second group fail) while the first group holds because $m>L/2$; there
$\Ext=\Nvis^\perp$ is still exact.  The dimension count of the brief
(\S\ref{sec:notes}~(N4)) is the finite-dimensional shadow of exactly this: $m=n/2$ at half filling
is the finite-dimensional generic-position condition.
\end{verdictok}

\begin{theorem}[Second, independent proof of (Q): duality]\label{thm:Qdual}
$\Nvis^\perp=i\overline{\mathcal K}=\overline{\Ext}$, where
$\mathcal K=\{\Pperp h'P:\ h'=h'^*\ \text{finite rank},\ h'=E_{I^c}h'E_{I^c}\}$.  In other words,
(Q) is precisely Haag (twisted) duality for the free chiral fermion net, read in the
one-particle--one-hole sector.  The form of duality used is
$\cF(I)'=Z\cF(I^c)Z^*$ with $I$ a \emph{bounded} interval and $I^c$ the union of two half-lines,
i.e.\ twisted duality for a disconnected complement, together with additivity
$\cF(I^c)=\cF(I^c_-)\vee\cF(I^c_+)$; for the free chiral fermion this is standard in the circle
picture, where $I^c$ is the complementary arc and the point at infinity has measure zero (Foit
1983, Thm.~1; PA Lemma~7.5 $\langle2\rangle1$ is written in PA's half-line configuration
$I=(-\infty,1)$, $I'=(1,\infty)$, and transfers by the M\"obius map of PA \S6.2 --- see
\S\ref{sec:corr} R6).  Theorem~\ref{thm:Q} needs none of this.
\end{theorem}
\lstep{1}{1}{Let $\cH^{(1,1)}\cong\cH$ be the one-particle--one-hole sector of Fock space, $E$ the
 orthogonal projection onto it, and $H_\cN:=\overline{\cN_{\rm sa}\Om}$ for a von Neumann algebra
 $\cN$.  Then $H_\cM\cap\cH^{(1,1)}=\overline\Nvis$ and
 $H_{\cM'}\cap\cH^{(1,1)}=\overline{\mathcal K}$, with $\cM=\cF(I)$.}
\lby{PA Lemma~7.5 (``Wick reduction'') $\langle1\rangle4$, whose proof uses twisted duality
 $\cM'=Z\cF(I')Z^*$ (Foit 1983, Thm.~1) only through the inclusion ``$\subseteq$'', and
 PA \S7.3 (its eq.~(35), the definition of $\mathcal L$ and $\mathcal K$): $\mathcal L=\Nvis$
 and $\mathcal K$ as above}
\lstep{1}{2}{For any closed real subspace $K$ of a complex Hilbert space,
 $K^{\perp_{\mathbb R}}=iK'$, where $K':=\{\xi:\operatorname{Im}\ip\xi\eta=0\ \forall\eta\in K\}$
 is the symplectic complement.}
\lby{$\Re\ip{i\zeta}\eta=\Re(-i\ip\zeta\eta)=\operatorname{Im}\ip\zeta\eta$, so
 $\xi=i\zeta\perp_{\mathbb R}K\iff\zeta\in K'$; pure algebra, no hypothesis}
\lstep{1}{3}{$\bigl(\overline{\cM_{\rm sa}\Om}\bigr)'=\overline{\cM'_{\rm sa}\Om}$.}
\lby{Tomita--Takesaki: for $\Om$ cyclic and separating (Reeh--Schlieder for $\cF(I)$, PA (C2)),
 the standard subspace $K=\overline{\cM_{\rm sa}\Om}$ has symplectic complement
 $\overline{\cM'_{\rm sa}\Om}$; ``$\supseteq$'' is elementary
 ($\omega(AB)=\overline{\omega(B^*A^*)}=\overline{\omega(BA)}=\overline{\omega(AB)}$ is real for
 $A\in\cM_{\rm sa}$, $B\in\cM'_{\rm sa}$), ``$\subseteq$'' is the standard-subspace form of
 Tomita's theorem (Rieffel--van Daele 1977, ``A bounded operator approach to Tomita--Takesaki
 theory'', Pacific J.~Math.~69, Thm.~2.4; also Longo's lecture notes ``Real Hilbert subspaces,
 modular theory, $SL(2,\mathbb R)$ and CFT'', Prop.~3.1.2).  \emph{Screenshots NOT OBTAINED}:
 neither source is in \texttt{refs/}}
\lstep{1}{4}{For $\xi\in\cH^{(1,1)}$: $\xi\perp_{\mathbb R}\Nvis\iff\xi\perp_{\mathbb R}H_\cM$.}
\lby{$E$ is a complex-linear orthogonal projection, so $\Re\ip\xi\eta=\Re\ip{E\xi}\eta
 =\Re\ip\xi{E\eta}$ for $\xi=E\xi$; and $EH_\cM\subseteq\overline\Nvis\subseteq H_\cM$
 ($\langle1\rangle1$ and PA Lemma~7.5)}
\lstep{1}{5}{$\Nvis^\perp=i\overline{\mathcal K}$.}
\lby{$\langle1\rangle4$, then $\langle1\rangle2$--$\langle1\rangle3$ give
 $\xi\perp_{\mathbb R}H_\cM\iff-i\xi\in\overline{\cM'_{\rm sa}\Om}=H_{\cM'}$; and $-i\xi\in
 \cH^{(1,1)}$, so $-i\xi\in H_{\cM'}\cap\cH^{(1,1)}=\overline{\mathcal K}$ by $\langle1\rangle1$}
\lstep{1}{6}{$i\mathcal K=\Ext_{\rm fr}$, hence $\Nvis^\perp=\overline{\Ext_{\rm fr}}
 =\overline\Ext$.}
\lby{$Z_e:=ih'$ runs over the finite-rank anti-self-adjoint operators supported in $I^c$ exactly
 when $h'$ runs over the finite-rank self-adjoint ones, and $\Pperp(ih')P=i\,\Pperp h'P$;
 then $\langle1\rangle5$ and the squeeze of Theorem~\ref{thm:Q} $\langle1\rangle1$}
\lqed{1}{7}\lby{$\langle1\rangle5$, $\langle1\rangle6$}

\section{Second-order control: from the linearised bound to the exact fidelity}\label{sec:second}

\begin{lemma}[One-parameter groups: the exact linear bound]\label{lem:group}
Let $(U_t)_{t\in\mathbb R}$ be a strongly continuous one-parameter unitary group on $H$ with
anti-self-adjoint generator $Z$, and assume $[P,Z]\in\Stwo$ (for unbounded $Z$: assume the
$\Stwo$-valued identity $U_t^*PU_t-P=\int_0^tU_\sigma^*[P,Z]U_\sigma\dd\sigma$, which holds for
every admissible field $Z=-D_\chi$ by \texttt{rate\_of\_remainder.tex} Lemma~6.1 $\langle1\rangle2$).  Then
\begin{equation}\label{eq:groupbound}
 \norm{P-U_t^*PU_t}_2\le\abs t\,\norm{[P,Z]}_2,\qquad
 \norm{\Pperp U_tP}_2\le\abs t\,\norm{\Pperp ZP}_2 .
\end{equation}
\end{lemma}
\lstep{1}{1}{For any unitary $U$ with $\Pperp UP\in\Stwo$: $\norm{P-U^*PU}_2^2=2\norm{\Pperp UP}_2^2$;
 and $\norm{[P,Z]}_2^2=2\norm{\Pperp ZP}_2^2$ for anti-self-adjoint $Z$.}
\lby{with $R:=U^*PU$ a projection, $\Tr((P-R)^2)=\Tr(P(\one-R))+\Tr(R(\one-P))
 =\norm{(\one-R)P}_2^2+\norm{(\one-P)R}_2^2=\norm{\Pperp UP}_2^2+\norm{P U\Pperp}_2^2$, and the two
 off-diagonal blocks of a unitary have the same singular values (CS decomposition,
 \texttt{rate\_of\_remainder.tex} Thm.~2.3 $\langle1\rangle4$); for the commutator,
 $[P,Z]=PZ\Pperp-\Pperp ZP$ (Lemma~\ref{lem:L1} $\langle1\rangle3$) is a sum of two operators with
 orthogonal ranges and orthogonal co-ranges, and $PZ\Pperp=-(\Pperp ZP)^*$}
\lstep{1}{2}{$\norm{P-U_t^*PU_t}_2\le\int_0^{\abs t}\norm{U_\sigma^*[P,Z]U_\sigma}_2\dd\sigma
 =\abs t\,\norm{[P,Z]}_2$.}
\lby{the hypothesis (for bounded $Z$: $\frac{\dd}{\dd t}U_t^*PU_t=U_t^*(Z^*P+PZ)U_t
 =U_t^*[P,Z]U_t$ in $\Stwo$, and a Bochner integral obeys the triangle inequality), and
 $\norm{V^*TV}_2=\norm T_2$ for unitary $V$}
\lqed{1}{3}\lby{$\langle1\rangle1$, $\langle1\rangle2$}

\begin{theorem}[Second-order control for a product family]\label{thm:UH}
Let $\chi$ be an admissible field (\texttt{rate\_of\_remainder.tex} Def.~2.1) with flow
$\kt_s$ and $\Wt_s:=W_{\kt_s}$; let $\Gs$ be bounded, anti-self-adjoint, $\Gs=E_D\Gs E_D$,
$\Pperp\Gs P\in\Stwo$; let $Z_e$ be bounded, anti-self-adjoint, $Z_e=E_{I^c}Z_eE_{I^c}$,
$\Pperp Z_eP\in\Stwo$; and let $\lambda\in\mathbb R$.  Put
\begin{equation}\label{eq:prodfam}
 U_s:=\Wt_s\,e^{-\lambda s\Gs}\,e^{sZ_e},\qquad
 L_0:=-D_\chi-\lambda\Gs+Z_e .
\end{equation}
Then $U_s\in\mathcal E(W_I^{(s,\lambda s)})$, where $W_I^{(s,\eps)}=\Wt_s e^{-\eps\Gs}|_{\mathfrak h_I}$
is the one-particle isometry of the rotated channel $\beta_{s,\eps}$ of S10 Lemma~6.1, and
\begin{equation}\label{eq:UHbound}
 -\log\mathrm F\bigl(\omega,\omega^{\beta_{s,\lambda s}}\bigr)
 \ \le\ \tfrac12\,s^2\,\bigl\|\Pperp L_0P\bigr\|_2^2\,\bigl(1+o(1)\bigr)\qquad(s\downarrow0).
\end{equation}
\end{theorem}
\lstep{1}{1}{$U_s$ is unitary, $U_s|_{\mathfrak h_I}=\Wt_se^{-\lambda s\Gs}|_{\mathfrak h_I}
 =W_I^{(s,\lambda s)}$, and $\Pperp U_sP\in\Stwo$; so $U_s\in\mathcal E(W_I^{(s,\lambda s)})$ and
 Prop.~\ref{prop:uhl} applies.}
\lby{$e^{sZ_e}$ fixes $\mathfrak h_I$ pointwise ($Z_e=E_{I^c}Z_eE_{I^c}$ gives $Z_eE_I=0$, so
 $e^{sZ_e}E_I=E_I$); S10 Lemma~6.1(1) for the channel properties of $\Wt_se^{-\eps\Gs}$ on
 $\mathfrak h_I$; $\Pperp U_sP\in\Stwo$ by $\langle1\rangle4$ below}
\lstep{1}{2}{$U_s$ is strongly differentiable with $\dot U_s=L_sU_s$,
 $L_s:=-D_\chi+\Wt_s(-\lambda\Gs)\Wt_s^*+\Wt_se^{-\lambda s\Gs}Z_ee^{\lambda s\Gs}\Wt_s^*$
 anti-self-adjoint, $L_0$ as in \eqref{eq:prodfam}, and
 $P-U_s^*PU_s=-\int_0^sU_\sigma^*[P,L_\sigma]U_\sigma\dd\sigma$ in $\Stwo$.}
\lby{the product rule for the three one-parameter groups (Leibniz on
 $\Wt_\sigma e^{-\lambda\sigma\Gs}e^{\sigma Z_e}$, with the two bounded generators moved to the
 left by conjugation), each factor being a strongly continuous group; then the argument of
 Lemma~\ref{lem:group} $\langle1\rangle2$ with the $\sigma$-dependent generator $L_\sigma$
 (\texttt{rate\_of\_remainder.tex} Lemma~6.1 $\langle1\rangle2$ for the unbounded factor $\Wt_\sigma$, the other two
 factors being norm-differentiable)}
\lstep{1}{3}{$\eta(\sigma):=\norm{[P,L_\sigma-L_0]}_2\to0$ as $\sigma\downarrow0$.}
\lproof
\lstep{2}{1}{For a bounded anti-self-adjoint $Y$ with $[P,Y]\in\Stwo$ and a unitary $V$:
 $[P,VYV^*]-[P,Y]=V[V^*PV-P,Y]V^*+\bigl(V[P,Y]V^*-[P,Y]\bigr)$.}
\lby{$[P,VYV^*]=V[V^*PV,Y]V^*$ (conjugate the commutator), then add and subtract
 $V[P,Y]V^*$}
\lstep{2}{2}{With $V=\Wt_\sigma$ resp.\ $V=\Wt_\sigma e^{-\lambda\sigma\Gs}$ and $Y=-\lambda\Gs$
 resp.\ $Y=Z_e$, both terms of $\langle2\rangle1$ tend to $0$ in $\Stwo$.}
\lby{first term: $\norm{[V^*PV-P,Y]}_2\le2\norm Y\,\norm{V^*PV-P}_2\to0$, because
 $\norm{P_\sigma-P}_2\le\sigma\norm{[P,D_\chi]}_2$ with $P_\sigma:=\Wt_\sigma^*P\Wt_\sigma$
 (Lemma~\ref{lem:group} applied to the group $\Wt_\sigma$), and in the second case
 $V^*PV-P=\bigl(e^{\lambda\sigma\Gs}(P_\sigma-P)e^{-\lambda\sigma\Gs}-(P_\sigma-P)\bigr)
 +\bigl(e^{\lambda\sigma\Gs}Pe^{-\lambda\sigma\Gs}-P\bigr)+(P_\sigma-P)$, of $\Stwo$-norm at most
 $3\norm{P_\sigma-P}_2+\lambda\sigma\norm{[P,\Gs]}_2\to0$ (Lemma~\ref{lem:group} again, for the
 group $e^{\lambda\sigma\Gs}$); second term: $\sigma\mapsto V_\sigma TV_\sigma^*$ is $\Stwo$-continuous at $\sigma=0$
 for $T\in\Stwo$ and a strongly continuous unitary family $V_\sigma$ (true for finite-rank $T$ by
 strong continuity, then by density and $\norm{V_\sigma TV_\sigma^*}_2=\norm T_2$)}
\lqed{2}{3}\lby{$\langle2\rangle1$, $\langle2\rangle2$ and linearity of $[P,\cdot]$}
\lstep{1}{4}{$\norm{\Pperp U_sP}_2\le s\norm{\Pperp L_0P}_2+\int_0^s\tfrac1{\sqrt2}\eta(\sigma)\dd\sigma
 =s\norm{\Pperp L_0P}_2\,(1+o(1))$.}
\lby{$\langle1\rangle2$, the triangle inequality for the Bochner integral,
 $\norm{[P,L_\sigma]}_2\le\norm{[P,L_0]}_2+\eta(\sigma)$, the identities of
 Lemma~\ref{lem:group} $\langle1\rangle1$, and $\int_0^s\eta=o(s)$ by $\langle1\rangle3$}
\lqed{1}{5}\lby{Prop.~\ref{prop:uhl} with $V=\Pperp U_sP$ and $\langle1\rangle4$:
 $-\log\mathrm F\le\tfrac12\norm V_2^2(1+O(\norm V_2^2))$}

\section{The upper half (UH), and the consequence for $\Erec$}\label{sec:main}

\begin{theorem}[(UH) for every isometric quasi-free channel]\label{thm:UHall}
Let $\beta_{s}$ be a family of isometric quasi-free channels realised by
$U_s=\Wt_s e^{-\lambda s\Gs}$ as in Theorem~\ref{thm:UH} (in particular $\lambda=0$ gives the
zero-collar compression).  Then
\begin{equation}\label{eq:UH}
 \limsup_{s\downarrow0}s^{-2}\bigl[-\log\mathrm F\bigl(\omega,\omega^{\beta_s}\bigr)\bigr]
 \ \le\ \tfrac12\dist\bigl(\Pperp(D_\chi+\lambda\Gs)P,\ \Ext\bigr)^2
 \ =\ \tfrac12\bigl\|\Pi\Pperp(D_\chi+\lambda\Gs)P\bigr\|_2^2 .
\end{equation}
Equivalently: the exact recovery error is, to second order, at most the \emph{optimised} Uhlmann
bound, and the optimisation over exterior extensions loses nothing.
\end{theorem}
\lstep{1}{1}{For every exterior $Z_e$ as in Theorem~\ref{thm:UH},
 $\limsup_{s\downarrow0}s^{-2}\bigl[-\log\mathrm F(\omega,\omega^{\beta_s})\bigr]
 \le\tfrac12\norm{\Pperp(D_\chi+\lambda\Gs-Z_e)P}_2^2$.}
\lby{Theorem~\ref{thm:UH}, \eqref{eq:UHbound}, and $\norm{\Pperp L_0P}_2=
 \norm{\Pperp(D_\chi+\lambda\Gs-Z_e)P}_2$ ($L_0=-(D_\chi+\lambda\Gs-Z_e)$)}
\lstep{1}{2}{The left side of $\langle1\rangle1$ is at most
 $\tfrac12\dist\bigl(\Pperp(D_\chi+\lambda\Gs)P,\Ext\bigr)^2
 =\tfrac12\norm{\Pi\Pperp(D_\chi+\lambda\Gs)P}_2^2$.}
\lby{the left side does not depend on $Z_e$, so one may take the infimum of the right sides;
 $-Z_e$ runs over the exterior generators together with $Z_e$; then
 Theorem~\ref{thm:reduction} $\langle1\rangle1$--$\langle1\rangle2$ and Theorem~\ref{thm:Q}
 ($\overline\Ext=\Nvis^\perp$ makes the second term of \eqref{eq:reduction} vanish)}
\lqed{1}{3}\lby{$\langle1\rangle1$, $\langle1\rangle2$}

\begin{remark}[The $g_Q$ form, and why it is not stated as the theorem]\label{rem:gQform}
The right-hand side of \eqref{eq:UH} is $s^{-2}g_Q(\delta(s))+o(1)$, with
$\delta(s)=E_I[P,s(D_\chi+\lambda\Gs)]E_I+o(s)$ the first-order defect (the sign is S10
eq.~(24); directly $\delta=E_I(P-U_s^*PU_s)E_I=-sE_I[P,L_0]E_I+O(s^2)$ with
$L_0=-(D_\chi+\lambda\Gs)$), \emph{whenever} S10 Lemma~6.1(2)--(3) applies: Lemma~\ref{lem:L2}
with $Z=s(D_\chi+\lambda\Gs)$ gives
$g_Q(E_I[P,Z]E_I)=\tfrac{s^2}2\norm{\Pi\Pperp(D_\chi+\lambda\Gs)P}_2^2$ exactly, and what is
needed in addition is $g_Q(\delta(s))=g_Q(E_I[P,Z]E_I)+o(s^2)$.  That last step is \emph{not} a
consequence of ``continuity of $g_Q$'' --- an earlier version of this document asserted so, and
it is false: $g_Q(\delta)=\tfrac18\sup_t(\Tr t\delta)^2/\Tr(tQt(\one-Q))$ is an unbounded
quadratic form because $\operatorname{spec}Q$ accumulates at $0$ and $1$ for an interval, which
is the documented order-of-limits effect (S10 \S7(d): uncapped
$g_Q(\delta_c)/\tfrac12s^2\norm u_2^2=2.8,12,45$ at $s=0.02,0.05,0.1$; S12 \S5.1: uncapped
$g_Q\sim10^{283}$ in the spectral model).  It is S10 Lemma~6.1(2)--(3) itself, with its weight
cap, that supplies the remainder control.  Nothing downstream uses the $g_Q$ form:
Theorem~\ref{thm:main} consumes \eqref{eq:UH} directly.  See \S\ref{sec:corr} R2.
\end{remark}

\begin{theorem}[Main]\label{thm:main}
Let
$\theta_{\rm fr}:=\sup\{\cos^2\angle(u,v_\Gs):\Gs\ \text{finite rank},\ v_\Gs\ne0\}$ over
finite-rank anti-self-adjoint $\Gs=E_D\Gs E_D$, with $u=\Pi M$, $M=\Pperp D_\chi P$,
$v_\Gs=\Pi(\Pperp\Gs P)$ as in S10 Def.~4.2 and Thm.~6.2.  Then, per chirality,
\begin{equation}\label{eq:main}
 \Erec\ \le\ (1-\theta_{\rm fr})\,f_2z^2\,(1+o(1))\qquad(z\to0),
\end{equation}
and $\theta_{\rm fr}>0$ if and only if $R_a:=\tfrac12(E_DuE_D-(E_DuE_D)^*)\ne0$, equivalently
(S10 Cor.~11.2: $E_DuE_D$ is purely anti-self-adjoint) if and only if $E_DuE_D\ne0$.
\end{theorem}
\lstep{1}{1}{For every finite-rank $\Gs$ and every $\lambda$,
 $\limsup_{z\to0}z^{-2}\Erec\le\tfrac12\norm{\Pi\Pperp(D_\chi+\lambda\Gs)P}_2^2$.}
\lby{$\Erec\le-\log\mathrm F(\omega,\omega^{\beta_s})$ for the particular channel $\beta_s$
 (definition of the infimum; $\beta_s$ is a legitimate recovery channel by S10 Lemma~6.1(1) and
 Theorem~\ref{thm:UH} $\langle1\rangle1$), Theorem~\ref{thm:UHall} $\langle1\rangle1$--$\langle1\rangle2$,
 and $\zeta=s=z$ in the half-line normal form (Note~5 Lemma~1.1: all quantities depend on
 $(a,L,s)$ only through the cross ratio); finite-rank $\Gs$ satisfies the hypothesis
 $\Pperp\Gs P\in\Stwo$ of Theorem~\ref{thm:UH}}
\lstep{1}{2}{$\min_\lambda\tfrac12\norm{\Pi\Pperp(D_\chi+\lambda\Gs)P}_2^2
 =\tfrac12\norm u_2^2\bigl(1-\cos^2\angle(u,v_\Gs)\bigr)$.}
\lby{$\Pi\Pperp(D_\chi+\lambda\Gs)P=u+\lambda v_\Gs$ (linearity of $\Pi$ and $u=\Pi M$,
 $M=\Pperp D_\chi P$), and $\min_\lambda\norm{u+\lambda v}^2=\norm u^2-
 (\Re\ip uv)^2/\norm v^2$; this is S10 Lemma~6.1(3) in the present notation}
\lstep{1}{3}{$\tfrac12\norm u_2^2=g_B/8=f_2$.}
\lby{S10 Prop.~4.3(2): $\norm u_2^2=g_B/4$, itself PA Prop.~7.6 (``the infimum is $g_B/4$'')
 together with \texttt{quadratic\_limit.tex} Thm.~4.2; and $f_2=g_B/8=1/(12\pi^2)$ for $c=1$.
 \emph{Note} (referee report \S10): this is where the main line does use net theory --- PA
 Prop.~7.6 rests on PA Lemma~7.5, hence on twisted duality --- although only the $X=I$ half of
 PA Lemma~7.5 is needed once Theorem~\ref{thm:Q} supplies
 $\dist(M,\Ext)^2=\norm{\Pi M}_2^2$ (the $X=I'$ half, which is the duality half, was what PA
 needed to \emph{identify} the infimum over extensions).  Theorem~\ref{thm:Q} and everything up
 to \eqref{eq:main} with $\norm u_2^2$ left unevaluated are duality-free; the \emph{number}
 $f_2$ is not}
\lstep{1}{4}{\eqref{eq:main}, and $\theta_{\rm fr}>0\iff R_a\ne0$.}
\lby{take the supremum over finite-rank $\Gs$ in $\langle1\rangle1$--$\langle1\rangle3$; for the
 equivalence, S10 Thm.~6.2 $\langle1\rangle1$ gives
 $\Re\ip u{v_\Gs}=\ip{R_a}\Gs_r$, so the supremum is $0$ iff $\ip{R_a}\Gs_r=0$ for all finite-rank
 anti-self-adjoint $\Gs=E_D\Gs E_D$, iff $R_a=0$ (the finite-rank anti-self-adjoint operators are
 $\Stwo$-dense and $R_a\in\Stwo$ is anti-self-adjoint, cf.\ Lemma~\ref{lem:L1} $\langle2\rangle2$)}
\lqed{1}{5}\lby{$\langle1\rangle4$}

\begin{remark}[The numerical rider: from $\theta_{\rm fr}$ to $0.38$]\label{rem:thetafr}
$\theta_{\rm fr}\le\theta:=\sup\{\cos^2\angle(u,v_\Gs):\Gs\ \text{bounded}\}$ by definition, so
the numerical statement $\theta=0.384\pm0.003$ (Phase~5, F3,
\texttt{numerics/optimality\_all/F3\_RESULTS.md}) does \emph{not} by itself give
$\Erec\le0.616f_2z^2$: what is needed is $\theta_{\rm fr}\ge0.38$.  Three observations supply it.
\begin{enumerate}[label=(\roman*),leftmargin=2.2em,itemsep=1pt]
\item $\theta_{\rm fr}=\theta_{\rm HS}:=\sup\{\cos^2\angle(u,v_\Gs):\Gs\in\Stwo\}$.  Indeed both
 $\Re\ip u{v_\Gs}=\Re\ip u{\Pperp\Gs P}$ and $\norm{v_\Gs}_2^2$ depend on $\Gs$ only through
 $w=\Pperp\Gs P$ and are $\Stwo$-continuous in $w$; $\Gs\mapsto\Pperp\Gs P$ is
 $\Stwo$-contractive; and the finite-rank anti-self-adjoint operators supported in $D$ are
 $\Stwo$-dense among the Hilbert--Schmidt ones.  Only $\theta_{\rm HS}=\theta$ (bounded $\Gs$
 with $\Pperp\Gs P\in\Stwo$ but $\Gs\notin\Stwo$) is left open, and it is irrelevant for a
 \emph{lower} bound on $\theta_{\rm fr}$.
\item There is an explicit admissible Hilbert--Schmidt direction: $\Gs=R_a$ (S10 Cor.~11.2
 $\langle1\rangle3$), giving
 $\theta_{\rm fr}\ge\norm{R_a}_2^4/(\norm u_2^2\ip{R_a}{\mathcal AR_a}_r)>0$ whenever
 $R_a\ne0$, numerically $\ge0.0143$.
\item The numerical optimisers are of this kind: S10 \S7(c) reports
 $\norm{\eps_*\Gs_*}_2/s=0.84$ (finite Hilbert--Schmidt norm), and every discretisation behind
 the value of $\theta$ is finite-dimensional, so its optimisers are finite rank.  This applies
 to the two frames reconciled in Phase~5 (F3): the modular Galerkin frame (S12 \S5.2) and the
 parameter-free modular Wiener--Hopf frame (\texttt{rigor/exact\_optimum\_tangent\_problem.tex}
 Prop.~6.2), which agree to $3\cdot10^{-4}$ after extrapolation and give
 $\theta=0.384\pm0.003$ (measured convergence order $h^{0.77}$); S10 \S7's $0.300\pm0.005$ is a
 taper-suppressed \emph{lower} bound, the ultraviolet smoothstep taper on the spectral
 derivative removing modes that carry $\approx20\%$ of $\theta$.  Since F3's extrapolations are
 themselves taken over finite meshes, $\theta_{\rm fr}\ge0.38$ is exactly as well supported
 numerically as $\theta=0.384\pm0.003$, and \eqref{eq:main} gives
 $\Erec\le0.616f_2z^2(1+o(1))$, i.e.\ $\kappa_{\rm opt}\le0.616\pm0.003$: \emph{the zero-collar
 compression is not optimal in exact fidelity}.  Whether non-isometric channels lower
 $\kappa_{\rm opt}$ further is the subject of Phase~5 and is not settled here.
\end{enumerate}
\end{remark}

\begin{corollary}[Sandwich]\label{cor:sandwich}
Per chirality, with $\mathcal D_z:=\min_\cF g_Q$ the value of the quasi-free second-order
programme,
\begin{equation}\label{eq:sandwich}
 \sup_{(t_z)}\ \liminf_{z\to0}\ \frac{\Delta_z(t_z)^2}{8V(t_z)z^2}
 \ \le\ \liminf_{z\to0}\frac{\Erec}{z^2}\ \le\ \limsup_{z\to0}\frac{\Erec}{z^2}
 \ \le\ \inf_{\lambda,\Gs}\ \tfrac12\bigl\|\Pi\Pperp(D_\chi+\lambda\Gs)P\bigr\|_2^2
 \ =\ (1-\theta_{\rm fr})f_2 ,
\end{equation}
the supremum over admissible test families.  The left-hand side equals $\mathcal D_z/z^2$ in the
limit along any family $t_z$ satisfying the conditions of S12 Thm.~4.1(2).  In particular, if
$\min_\cF g_Q$ is attained (in the limit) at an isometric point and is certified by such a
family, then $\Erec=\mathcal D_z(1+o(1))$ exactly.
\end{corollary}
\lby{the upper half is Theorem~\ref{thm:UHall} together with Theorem~\ref{thm:main}
 $\langle1\rangle2$--$\langle1\rangle3$.  The lower half is S12 Thm.~4.1, in its unconditional
 dual form (S12 eq.~(10)): $\liminf z^{-2}\Erec\ge\sup_{(t_z)}\liminf
 z^{-2}\Delta_z(t_z)^2/(8V(t_z))$, which comes from the exact single-observable bound S12
 Prop.~2.5.  The identification of that supremum with $\mathcal D_z=\min_\cF g_Q$ (S12
 Thm.~4.1(2), via Sion's minimax theorem) is \emph{conditional}: it requires a test family with
 $\norm{t_z}_1^2/V(t_z)$ bounded and $\Delta_z(t_z)\norm{t_z}_1/V(t_z)\to0$.  An earlier version
 of this corollary quoted ``S12 Prop.~2.4 and Thm.~4.1'' for an unconditional
 $\Erec\ge\mathcal D_z(1-o(1))$; see \S\ref{sec:corr} R5}

\section{Corroboration from the numerical track (E2)}\label{sec:E2}

Track E2 ran the Phase-4 numerics in S10's spectral NS-circle model; the results
(\texttt{rigor/exterior\_extension\_numerics.tex}, cards \texttt{cft\_cmi:e2-*}) match the
theorems above point by point.
\begin{enumerate}[leftmargin=1.8em,itemsep=2pt]
\item \emph{Theorem~\ref{thm:Q} in the model.}  $\Ext=\Nvis^\perp$ \emph{exactly} at every $L$,
 with $\dim_{\mathbb R}\Nvis=m^2-2z^2$, $\dim_{\mathbb R}\Nvis^\perp=(L-m)^2=\dim_{\mathbb R}\Ext$
 ($z=m-L/2$), and all principal angles between $\Ext$ and $\Nvis^\perp$ equal to $0$ to $10$
 digits at $L=32,48,64$ (card \texttt{e2-ext-equals-nperp-circle}).  Note that the finite model is
 in generic position in the sense of Lemma~\ref{lem:generic} only up to the $2z^2$ exact
 $\{0,1\}$-eigenvalues of $Q$, which E2 accounts for explicitly; this is the finite-dimensional
 shadow of \S\ref{sec:notes}~(N4).
\item \emph{Theorem~\ref{thm:reduction} and Corollary~\ref{cor:UHneeds}.}
 $\dist((\one-\Pi)M,\Ext)^2$ and $\dist((\one-\Pi)v_{\Gs_*},\Ext)^2$ vanish to $55$--$70$ digits
 in \texttt{mpmath} (card \texttt{e2-mpmath-distance-exactly-zero}); the optimised Uhlmann bound
 gives $\mathrm{UB}_c/(\tfrac12s^2\norm u_2^2)=1.0000000$ and
 $\mathrm{UB}_{\rm rot}/\mathrm{UB}_c=1-\theta$ to $4$--$5$ digits at $L=128,\dots,384$, and is
 exactly quadratic in $s$ (card \texttt{e2-optimised-uhlmann}).  The factor $4.3$ of the earlier
 Gaussian-taper extension (\texttt{numerics/optimality\_all/uhlmann\_rotated.py}) was entirely an
 artefact of a bad exterior extension.
\item \emph{The infimum is essentially attained, with bounded norm.}  The optimal exterior
 generator has $\norm{Z_e}_F/s=3.03,3.13,3.18,3.22$ at $L=128,192,256,384$ --- bounded uniformly
 in $L$.  No uniformity in $Z_e$ is claimed or needed: Theorem~\ref{thm:UHall}
 $\langle1\rangle1$--$\langle1\rangle2$ takes the $\limsup_{s\downarrow0}$ at fixed $Z_e$ and only
 then the infimum over $Z_e$, which is the hypothesis-free order.  (An earlier version of this
 item asserted that the $o(1)$ of Theorem~\ref{thm:UH} is uniform on norm-bounded sets of
 exterior generators; that is false --- the modulus of continuity of
 $\sigma\mapsto V_\sigma[P,Z_e]V_\sigma^*$ at $\sigma=0$ depends on the operator $[P,Z_e]$ and not
 only on its norm --- and it is not used; see \S\ref{sec:corr} R9.)  What E2's boundedness does
 show is that the infimum is effectively attained.
\item \emph{Conditioning.}  E2 also reports that a growing fraction of $\Ext$-directions requires
 an exponentially large $\norm{Z_e}_F$ as $L$ grows.  That is the finite-$L$ signature of
 ``$\Ext$ dense but not closed'', consistent with Theorem~\ref{thm:Q} being a \emph{closure}
 statement (verdict after Theorem~\ref{thm:Q}); the directions we actually need,
 $(\one-\Pi)M$ and $(\one-\Pi)v_{\Gs_*}$, are not among them.
\end{enumerate}

\section{Verdict}\label{sec:verdict}

\begin{verdictok}[\;What is proved here]
\begin{enumerate}[leftmargin=1.6em,itemsep=1pt]
\item \textbf{(Q) is true}: $\overline{\Ext}=\Nvis^\perp$ in $\Stwo(PH,\Pperp H)$
 (Theorem~\ref{thm:Q}), with finite-rank exterior generators already dense.  Two independent
 proofs: an invariant two-projection computation (inputs: generic position of $(P,E_I)$, i.e.\
 Hardy-space uniqueness, Lemma~\ref{lem:generic}; and compactness of the test element), and
 Tomita--Takesaki duality in the $(1,1)$ sector (Theorem~\ref{thm:Qdual}), which shows that (Q)
 \emph{is} Haag duality read in that sector.  No numerical input.
\item \textbf{The reduction}: for any anti-self-adjoint $Z_0$ with $\Pperp Z_0P\in\Stwo$,
 $\inf_{Z_e}\tfrac12\norm{\Pperp(Z_0+Z_e)P}_2^2=\tfrac12\dist(\Pperp Z_0P,\Ext)^2
 =g_Q(E_I[P,Z_0]E_I)$ (Theorems~\ref{thm:reduction} and \ref{thm:Q}), and the value depends on the
 channel only through its first-order defect (Corollary~\ref{cor:extindep}).  The correct
 bookkeeping is the prescribed column / free exterior block (Lemma~\ref{lem:column}), not a
 block-diagonal generator.
\item \textbf{(UH)}, in the $g_Q$-free form \eqref{eq:UH}:
 $\limsup s^{-2}[-\log\mathrm F(\omega,\omega^{\beta_s})]\le
 \tfrac12\norm{\Pi\Pperp(D_\chi+\lambda\Gs)P}_2^2$ for the rotated family and for the
 compression (Theorems~\ref{thm:UH}, \ref{thm:UHall}); this equals $s^{-2}g_Q(\delta(s))+o(1)$
 whenever S10 Lemma~6.1(2)--(3) applies (Remark~\ref{rem:gQform}).  Via the exact
 Uhlmann bound (Prop.~\ref{prop:uhl}) and the one-parameter-group estimate
 (Lemma~\ref{lem:group}).  This is the missing upper half of the exact-fidelity question.
\item \textbf{Consequence}: $\Erec\le(1-\theta_{\rm fr})f_2z^2(1+o(1))$ with
 $\theta_{\rm fr}>0\iff E_DuE_D\ne(E_DuE_D)^*$ (Theorem~\ref{thm:main}), and the sandwich
 $\min_\cF g_Q(1-o(1))\le\Erec\le\inf_{\rm isom}g_Q(1+o(1))$ (Corollary~\ref{cor:sandwich}).
\end{enumerate}
\end{verdictok}

\begin{verdictgap}[\;What is \emph{not} proved]
\begin{enumerate}[leftmargin=1.6em,itemsep=1pt]
\item \textbf{$\theta_{\rm fr}>0$ is still numerical.}  Nothing here proves
 $E_DuE_D\ne0$; that remains S10 \S7 and S12 \S5.1--5.2, reconciled in Phase~5 (F3) to
 $\theta=0.384\pm0.003$ from two independent discretisations (modular Galerkin and modular
 Wiener--Hopf), S10's $0.300\pm0.005$ being a taper-suppressed lower bound and E2's
 $\mathrm{UB}_{\rm rot}/\mathrm{UB}_c=0.708\ldots0.711$ at $L=128\ldots384$ reconfirming the
 same tapered circle-model value.  The inequality \eqref{eq:main} with
 $\theta_{\rm fr}$, and the equivalence $\theta_{\rm fr}>0\iff E_DuE_D\ne0$, are proved; the
 \emph{numerical} statements are $\theta_{\rm fr}\ge0.38$ (Remark~\ref{rem:thetafr}) and hence
 $\kappa_{\rm opt}\le0.616\pm0.003$.
\item \textbf{The rotation class.}  Theorem~\ref{thm:UH} assumes $\Pperp\Gs P\in\Stwo$, which is
 \emph{not} automatic for bounded $\Gs=E_D\Gs E_D$ (for $\Gs=i\one_D$, $\Pperp\Gs P$ is not
 Hilbert--Schmidt: that is the logarithmic divergence of the sharp-interval entanglement).
 Remark~\ref{rem:thetafr}(i) closes the gap between $\theta_{\rm fr}$ and $\theta_{\rm HS}$;
 what is \emph{not} proved is $\theta_{\rm HS}=\theta$, the supremum over all bounded $\Gs$ with
 $\Pperp\Gs P\in\Stwo$.  This is harmless: only a lower bound on $\theta_{\rm fr}$ is used.
\item \textbf{No rate.}  The $o(1)$ of Theorem~\ref{thm:UH} comes from an $\Stwo$-continuity, not
 from a Lipschitz estimate, so no explicit remainder is given; combined with
 \texttt{rate\_of\_remainder.tex} Thm.~2.6 one expects a relative $O(s^2)$, but this is not proved
 for the product family.
\item \textbf{Cited, not verified here}: Hardy-space uniqueness (Rudin, Thms.~17.18 and 19.2,
 textbook, \emph{screenshot NOT OBTAINED}); for the second proof only, twisted duality (Foit 1983,
 via PA Lemma~7.5) and the standard-subspace form of Tomita's theorem (Rieffel--van~Daele 1977,
 \emph{screenshot NOT OBTAINED}).  The first proof (Theorem~\ref{thm:Q}) uses none of these ---
 but the \emph{number} $f_2=\tfrac12\norm u_2^2$ in \eqref{eq:main} does, through S10
 Prop.~4.3(2) and PA Prop.~7.6 (Theorem~\ref{thm:main} $\langle1\rangle3$): the structural chain
 is duality-free, the constant is not.
\item \textbf{Non-isometric channels}: $\min_\cF g_Q$ versus $\inf_{\rm isom}g_Q$ is untouched
 (Corollary~\ref{cor:sandwich}); the exact value of $\mathcal D_z$ remains open.
\end{enumerate}
\end{verdictgap}

\begin{verdictbreak}[\;Correction to the Phase-4 brief]
The brief asserts that an admissible extension $\Wt$ ``preserves $L^2(I^c)$ as well'' and writes
$Z=Z_I\oplus Z_e$ block diagonally.  Both are false: $W_I$ is an isometry of $\mathfrak h_I$
\emph{onto the proper subspace} $\mathfrak h_{AB}$, so $\Wt\mathfrak h_{I^c}\supsetneq
\mathfrak h_{I^c}$, and no admissible generator is block diagonal.  The statements survive with
$Z_I$ read as the prescribed \emph{column} $ZE_I$ (Lemma~\ref{lem:column}); reading it as the
truncation $E_ID_\chi E_I$ produces a visible defect that diverges like $L^2$
(Remark~\ref{rem:trunc}).  Nothing else in the brief's logic is affected.
\end{verdictbreak}

\section{Corrections after the referee report}\label{sec:corr}

The referee report \texttt{rigor/referee\_exact\_fidelity\_upper\_half.tex} (14 pp.) found no
proof-breaking issue --- ``the document stands, with repairs'' --- and requested R1--R10 of its
\S13.  All are applied \emph{in place} above; this section records what changed, so that a reader
of an earlier copy can see the differences.

\begin{enumerate}[leftmargin=2.6em,itemsep=2pt]
\item[\textbf{R1}] \emph{From $\theta$ to $\theta_{\rm fr}$.}  Theorem~\ref{thm:main} now states
 only \eqref{eq:main} with $\theta_{\rm fr}$ (supremum over finite-rank $\Gs$ with $v_\Gs\ne0$)
 and the equivalence $\theta_{\rm fr}>0\iff E_DuE_D\ne0$, both proved.  The numerical conclusion
 $\kappa_{\rm opt}\le1-\theta_{\rm fr}$ (then quoted as $0.71$, now $0.616\pm0.003$; see the
 addendum at the end of this section) moved to Remark~\ref{rem:thetafr}, which supplies the
 missing bridge:
 (i) $\theta_{\rm fr}=\theta_{\rm HS}$ by $\Stwo$-density, (ii) the explicit Hilbert--Schmidt
 direction $\Gs=R_a$ of S10 Cor.~11.2, (iii) the discretisations behind the numerical $\theta$
 are finite-dimensional, hence their optimisers are finite rank.  Only $\theta_{\rm HS}=\theta$
 remains open, and it is not needed for a lower bound on $\theta_{\rm fr}$.
\item[\textbf{R2}] \emph{The $g_Q$ form of (UH).}  Theorem~\ref{thm:UHall} is now stated in the
 proved, $g_Q$-free form \eqref{eq:UH}.  The identification with $s^{-2}g_Q(\delta(s))$ is
 Remark~\ref{rem:gQform}, which replaces the false justification ``continuity of the positive
 quadratic form $g_Q$'' by the correct one (S10 Lemma~6.1(2)--(3) with its weight cap) and
 records why continuity fails (S10 \S7(d); S12 \S5.1).  Corollary~\ref{cor:sandwich} and the
 verdict are rephrased accordingly; Theorem~\ref{thm:main} was always consuming
 \eqref{eq:UH}, not the $g_Q$ form, so nothing else changed.
\item[\textbf{R3}] \emph{Compactness is sufficient, not necessary.}  The verdict box after
 Theorem~\ref{thm:Q} no longer claims that $\Gamma'\pi\Gamma'=\pi$ has nonzero bounded solutions
 (it does not: approximate eigenvectors at $\pm1$ tend weakly to $0$).  It now records the
 referee's two-sided estimate $\abs{\ip\xi{\pi\eta}}\le\norm{\Gamma'^{\,n}\xi}\norm\pi
 \norm{\Gamma'^{\,n}\eta}\to0$, which kills every bounded solution, so step
 $\langle1\rangle4$ holds with ``compact'' deleted; and it gives the correct reason why only a
 closure is claimed (E2 \S2.3: $Z_e\mapsto\Pperp Z_eP$ has unbounded inverse on $\Nvis^\perp$).
\item[\textbf{R4}] \emph{Lemma~\ref{lem:column}(c).}  Restated as the differentiated torsor
 statement (right translation by $e^{\tau Z_e}$), with the caveat that beyond first order the
 four blocks of $Z$ are not independent.  Theorem~\ref{thm:UH} uses an exact right translation,
 so nothing downstream changes.
\item[\textbf{R5}] \emph{The lower half of the sandwich.}  Corollary~\ref{cor:sandwich} now
 states the unconditional dual form (S12 eq.~(10), from S12 Prop.~2.5) and says explicitly that
 the identification of that supremum with $\mathcal D_z=\min_\cF g_Q$ (S12 Thm.~4.1(2)) is
 conditional on an admissible test family.  The locator ``S12 Prop.~2.4'' (nonexistent) is
 corrected to Prop.~2.5.
\item[\textbf{R6}] \emph{Geometry transfer.}  Theorem~\ref{thm:Qdual} now names the form of
 twisted duality used for a \emph{bounded} interval (disconnected complement, plus additivity;
 circle picture), and says that PA Lemma~7.5 is written in PA's half-line configuration.
 Theorem~\ref{thm:Q} is unaffected: Lemma~\ref{lem:generic} holds for any $I$ with $I$ and $I^c$
 of positive measure.
\item[\textbf{R7}] \emph{Scope of ``duality-free''.}  Theorem~\ref{thm:main} $\langle1\rangle3$
 and the verdict now say that the \emph{value} $\tfrac12\norm u_2^2=f_2$ reaches back to PA
 Prop.~7.6 and hence to twisted duality, while the structural chain
 (Theorems~\ref{thm:Q}, \ref{thm:reduction}, \ref{thm:UH}, \ref{thm:UHall}) is duality-free.
\item[\textbf{R8}] \emph{Locators.}  The 15 references listed in the report's \S12 are
 renumbered: PA Cor.~1.9$\to$2.7, PA \S2$\to$\S3, PA Thm.~3.1$\to$4.1, PA Lemma~2.4$\to$3.2,
 PA Lemma~5.5$\to$7.5, PA \S5.3 eq.~(5.5)$\to$\S7.3 eq.~(35), PA Prop.~5.9$\to$7.6,
 RR Lemma~2.2$\to$6.1 $\langle1\rangle2$, RR Thm.~2.5 $\langle1\rangle4\to$2.3
 $\langle1\rangle4$, RR Thm.~2.6 $\langle1\rangle2\to$2.5 $\langle1\rangle2$,
 S10 eq.~(2.1)$\to$(1), S10 \S9$\to$\S7, S12 \S5.2--5.3$\to$\S5.1--5.2,
 S12 Prop.~2.4$\to$2.5.  ``Cousin--Sylvester'' is corrected to ``cosine--sine'', and the sigil
 S12 $=$ \texttt{optimality\_all\_channels.tex} is defined at its first use in (E3).
\item[\textbf{R9}] \emph{Two asides.}  The sign of $\delta(s)$ now matches S10 eq.~(24)
 (Remark~\ref{rem:gQform}); and the false uniformity claim in \S\ref{sec:E2}, item~3 (that the
 $o(1)$ of Theorem~\ref{thm:UH} is uniform on norm-bounded sets of exterior generators) is
 deleted --- the correct order of limits, $\limsup_{s\to0}$ at fixed $Z_e$ and only then
 $\inf_{Z_e}$, is what Theorem~\ref{thm:UHall} uses.
\item[\textbf{R10}] \emph{Abstract.}  $\Nvis$ now carries the finite-rank restriction,
 $\theta_{\rm fr}$ replaces $\theta$ with the proviso $v_\Gs\ne0$, and the description of the
 first proof is corrected from ``Halmos two-projection normal form'' to ``invariant
 two-projection computation'' (the Halmos form appears only in the working notes \S\ref{sec:notes}
 (N2), where the argument was found).
\end{enumerate}

\begin{verdictok}
After R1--R10 the load-bearing chain is: Theorem~\ref{thm:Q} ($\overline\Ext=\Nvis^\perp$,
proved, no compactness, no duality, no numerics) $\Rightarrow$ Theorem~\ref{thm:reduction} and
Corollary~\ref{cor:extindep} $\Rightarrow$ Theorem~\ref{thm:UH} and \eqref{eq:UH}
$\Rightarrow$ \eqref{eq:main}, $\Erec\le(1-\theta_{\rm fr})f_2z^2(1+o(1))$ with
$\theta_{\rm fr}>0\iff E_DuE_D\ne0$.  The only numerical inputs are the \emph{value}
$\theta_{\rm fr}\ge0.38$ ($\theta=0.384\pm0.003$, F3) and, through $f_2$, the constant $g_B=2/(3\pi^2)$ of
\texttt{quadratic\_limit.tex} Thm.~4.2 (which is proved, not numerical).  No statement of the
document is now known to be false.
\end{verdictok}

\paragraph{Addendum (2026-09-15): the value of $\theta$.}  After the referee report, Phase~5
track~F3 (\texttt{numerics/optimality\_all/F3\_RESULTS.md}) reconciled the two numerical frames
for $\theta$.  The outcome: $\theta=0.384\pm0.003$, geometry independent, measured convergence
order $h^{0.77}$, obtained from two independent discretisations --- the modular Galerkin frame
(S12 \S5.2) and the parameter-free modular Wiener--Hopf frame
(\texttt{rigor/exact\_optimum\_tangent\_problem.tex} Prop.~6.2) --- which agree to
$3\cdot10^{-4}$ after extrapolation.  S10 \S7's $\theta=0.300\pm0.005$, quoted throughout the
earlier version of this document and in the referee report, is a \emph{taper-suppressed lower
bound}: the smoothstep ultraviolet taper S10 puts on the spectral derivative removes modes
carrying $\approx20\%$ of $\theta$, there is no taper-free lattice limit (the periodic grid's
second Fermi point), and $\theta_{\rm circ}$ rises linearly in $\log L$, so the $1/L$ Richardson
extrapolation used the wrong variable.  Running the modular-Galerkin functional on the circle
model's own mesh reproduces the Galerkin values.  Accordingly the numerical riders above now read
$\theta_{\rm fr}\ge0.38$ and $\kappa_{\rm opt}\le0.616\pm0.003$ in place of
$\theta_{\rm fr}\ge0.30$ and $\kappa_{\rm opt}\le0.71$; whether non-isometric channels lower
$\kappa_{\rm opt}$ further is the subject of Phase~5.  Nothing proved in this document changes:
all statements are inequalities in $\theta_{\rm fr}$, and only the numerical input moved.  S10
stays as the historical record; the erratum is \texttt{rigor/errata\_log.md}.

\end{document}
